\section{The class of 79}\label{sec:79}

In this section $f=93$ and $F=F_{93}$. The rule is $f(a,b,c)=\neg c\vee(\neg a\wedge b)$, with values $1,0,1,1,1,0,1,0$ on $000,\dots,111$. Reflection $f(a,b,c)\mapsto f(c,b,a)$ takes $93$ to $79$, complementation $f(a,b,c)\mapsto f(1-a,1-b,1-c)$ takes $93$ to $186$ and $79$ to $242$, so by Section~\ref{sec:conj} it suffices to treat $93$. The rule does not satisfy \eqref{eq:K}, since $f(0,0,0)=1$, and the constant configurations are not a quiet background; their role is taken by a configuration of period $2$ in space and $6$ in time, the \emph{vacuum}. We use the following values of $f$:
\begin{equation}\label{c79:eq:f}
f(a,b,0)=1,\qquad f(a,b,1)=\neg a\wedge b,\qquad f(a,1,1)=\neg a,\qquad f(0,1,c)=1,
\end{equation}
and $f(0,b,c)=f(1,b,c)$ unless $(b,c)=(1,1)$. On the constant configurations, $F$ acts by $(0,0)^{\Z}\mapsto(0,1)^{\Z}\mapsto(1,0)^{\Z}\mapsto(0,0)^{\Z}$ and fixes $(1,1)^{\Z}$, since $f(0,0,0)=1$ and $f(1,1,1)=0$.

The proof follows the plan of Section~\ref{cl:sec:110}, with three changes. The edge is the last cell at which a configuration differs from a vacuum, and a move changes the phase of the vacuum to the right of the edge rather than the edge cell. The second and third columns are read off by moves at a few times next to each time, using the period-$6$ structure of the edge column. And the controllability theorem starts at $k=4$, from a statement at level $3$ that holds up to five exceptional configurations and is proved from a table of the moves at the times $|m|\le3$.

\subsection{The vacuum and the edge}\label{c79:sec:setup}

Let $V\colon\Z^2\to\{0,1\}$ be given by
\[
V(t,i)=\begin{cases}[\,i\text{ odd}\,]&t\equiv0,1\pmod 6,\\ 0&t\equiv2,5\pmod 6,\\ [\,i\text{ even}\,]&t\equiv3,4\pmod 6,\end{cases}
\]
and for $\omega\in\Z/6$ let $V_\omega$ be the configuration with cells $(V_\omega)_i=(V(\omega-1,i),V(\omega,i))$, the \emph{vacuum of phase $\omega$}. Its cells at an even and an odd site are $(0,0),(0,1)$ for $\omega=0$; $(0,0),(1,1)$ for $\omega=1$; $(0,0),(1,0)$ for $\omega=2$; $(0,1),(0,0)$ for $\omega=3$; $(1,1),(0,0)$ for $\omega=4$; and $(1,0),(0,0)$ for $\omega=5$.

\begin{lemma}[vacuum]\label{c79:lem:vac}
\begin{enumerate}
\item[(a)] $V$ satisfies the spacetime relation \eqref{eq:rel}, so $V$ is the diagram of $V_0$ and $F^tV_0=V_t$. Moreover $\sh V_\omega=V_{\omega+3}$, $\sh^2V_\omega=V_\omega$, $F^3V_\omega=\sh V_\omega$ and $\rho V_\omega=V_{2-\omega}$.
\item[(b)] Two vacua that agree at two adjacent sites are equal, and $V_\omega$ and $V_{\omega+3}$ differ at every site.
\item[(c)] \emph{(Left-blindness.)} $(V(t,i),V(t,i+1))\ne(1,1)$ for all $t,i$. Consequently, if $d_{i}$ and $d_{i+1}$ agree with $V_\omega$, then $F(d)_i$ and $F^{-1}(d)_i$ do not depend on $d_{i-1}$.
\end{enumerate}
\end{lemma}

\begin{proof}
(a) Row $t$ of $V$ is $[i\text{ odd}]$, $[i\text{ odd}]$, $0$, $[i\text{ even}]$, $[i\text{ even}]$, $0$ for $t\equiv0,\dots,5$. On a row $[i\text{ odd}]$ or $[i\text{ even}]$ the triple at $i$ is $(0,1,0)$ or $(1,0,1)$, and $f(0,1,0)=1$, $f(1,0,1)=0$, so the right side of \eqref{eq:rel} is the row itself; on a zero row it is $f(0,0,0)=1$. In each of the six cases the left side $V(t+1,i)\oplus V(t-1,i)$ is the same. The row of $V$ at time $t+3$ is that at time $t$ shifted by one site, which gives $\sh V_\omega=V_{\omega+3}$ and $F^3V_\omega=V_{\omega+3}$; the period $2$ is clear. By \eqref{eq:rho}, $\rho V_\omega$ has diagram $V(\omega-1-t,i)$, and $V(1-t,i)=V(t,i)$ for all $t$, so $V(\omega-1-t,i)=V(t+2-\omega,i)$. (b) is read off the list of cells above. (c) No row of $V$ has two adjacent ones, and $f(a,b,c)$ depends on $a$ only if $(b,c)=(1,1)$. The $y$-row and the $x$-row of $V_\omega$ are rows of $V$, and $F(d)_i$ reads $a=y_{i-1}$ with $(b,c)=(y_i,y_{i+1})$, while $F^{-1}(d)_i$ reads $x_{i-1}$ with $(x_i,x_{i+1})$.
\end{proof}

For $r\in\Z$ let $\Omega_r$ be the set of configurations $c$ such that, for some $\omega$, $c_i=(V_\omega)_i$ for all $i>r$ and $c_r\ne(V_\omega)_r$; by Lemma~\ref{c79:lem:vac}(b), $\omega$ is unique, and we call it the \emph{phase} $\omega(c)$ and $r$ the \emph{edge} of $c$. Put $\vartheta(c)=\omega(c)+3r\in\Z/6$; by Lemma~\ref{c79:lem:vac}(a), $\sh^rc\in\Omega_0$ has phase $\vartheta(c)$. For $c\in\Omega_r$ write $z_t=Y_c(t,r)$ for the edge column. If $c$ agrees with $V_\omega$ on $(r,\infty)$ and $\Lambda$ agrees with $V_\omega$ on $(-\infty,p)$ with $p\ge r+1$, let $c\sqcup_r\Lambda$ be the configuration equal to $c$ on $(-\infty,r]$ and to $\Lambda$ on $(r,\infty)$; it equals $c\sqcup_{r'}\Lambda$ for every $r'\in[r,p-1]$. The same notation, with the same meaning, is used when $c$ agrees with a vacuum $V_\omega$ on $(r,\infty)$ and $\Lambda$ is a vacuum $V_{\omega'}$ of a possibly different phase, as in Definition~\ref{c79:def:moves}; Lemmas~\ref{c79:lem:glue}(a) and~\ref{c79:lem:hglue} assume a common phase.

\begin{lemma}[gluing and the edge]\label{c79:lem:glue}
\begin{enumerate}
\item[(a)] Let $c$ agree with $V_\omega$ on $(r,\infty)$ and $\Lambda$ with $V_\omega$ on $(-\infty,p)$, where $p\ge r+3$. Then $F^{\pm1}(c\sqcup_r\Lambda)=F^{\pm1}c\sqcup_rF^{\pm1}\Lambda$, where $F^{\pm1}c$ agrees with $V_{\omega\pm1}$ on $(r,\infty)$ and $F^{\pm1}\Lambda$ agrees with $V_{\omega\pm1}$ on $(-\infty,p-1)$.
\item[(b)] If $c\in\Omega_r$, then $F^tc\in\Omega_r$ with phase $\omega(c)+t$, for every $t\in\Z$. The columns of $Y_c$ to the right of $r$ are those of $V$ shifted by $\omega(c)$ in time.
\item[(c)] For $c\in\Omega_r$ and $t\in\Z$, with $s_t=V(t+\omega(c),r+1)$,
\[
z_{t+1}=z_{t-1}\oplus f\bigl(Y_x(t,r-1),z_t,s_t\bigr),
\]
and $s_t=1$ if and only if $t+\vartheta(c)\equiv0,1\pmod 6$. So $z_{t+1}=\neg z_{t-1}$ if $t+\vartheta(c)\equiv2,3,4,5$, and $z_{t+1}=z_{t-1}\oplus(z_t\wedge\neg Y_c(t,r-1))$ otherwise. The cell of $V_\omega$ at $r$ at time $t$ has $y$-bit $1$ if and only if $t+\omega+3r\equiv3,4\pmod 6$.
\end{enumerate}
\end{lemma}

\begin{proof}
(a) Put $\hat c=c\sqcup_r\Lambda$, which equals $c$ on $(-\infty,p-1]$ and $\Lambda$ on $(r,\infty)$. For $i\le p-2$ the cell $F^{\pm1}({\hat c})_i$ reads $\hat c$ on $[i-1,i+1]\subseteq(-\infty,p-1]$, so it is $F^{\pm1}(c)_i$; for $i\ge r+2$ it reads only cells of $\Lambda$; and for $i=r+1$ it reads ${\hat c}_{r+1}=\Lambda_{r+1}$ and ${\hat c}_{r+2}=\Lambda_{r+2}$, which agree with $V_\omega$, so by Lemma~\ref{c79:lem:vac}(c) it does not depend on ${\hat c}_r$ and equals $F^{\pm1}(\Lambda)_{r+1}$. As $p-2\ge r+1$, this is the identity. Applying the same argument to $c$ glued to $V_\omega$, and to $V_\omega$ glued to $\Lambda$, gives the statements about $F^{\pm1}c$ and $F^{\pm1}\Lambda$, since $F^{\pm1}V_\omega=V_{\omega\pm1}$.

(b) By (a), $F^{\pm1}c$ agrees with $V_{\omega\pm1}$ on $(r,\infty)$. If $F(c)_r$ were $(V_{\omega+1})_r$, then $F^{-1}(Fc)_r$ would read the cells $r,r+1$ of $Fc$, which agree with $V_{\omega+1}$, so by Lemma~\ref{c79:lem:vac}(c) it would equal $(V_\omega)_r$, contradicting $c\in\Omega_r$. The same argument applies to $F^{-1}$; induct on $|t|$.

(c) This is \eqref{eq:rel} at $(t,r)$, with $Y_c(t,r+1)=V(t+\omega,r+1)$ by (b), and \eqref{c79:eq:f}. If $r$ is even, then $\vartheta\equiv\omega$, and $V(t+\omega,r+1)=1$ exactly when $t+\omega\equiv0,1$; if $r$ is odd, then $\vartheta\equiv\omega+3$, and $V(t+\omega,r+1)=1$ exactly when $t+\omega\equiv3,4$. The last statement is the same computation at column $r$.
\end{proof}

\subsection{Kinks and normalisation}\label{c79:sec:kinks}

For $s=\pm1$ let $\mathrm{Fix}_s$ be the set of fixed points of $\sh^{-s}F^3$. As in Lemma~\ref{cl:lem:Z2}, these are coded by sequences. Let $\Phi$ be the set of $\varphi\in\{0,1\}^{\Z}$ with
\begin{equation}\label{c79:eq:Phi}
\varphi(n-1)\oplus\varphi(n+1)=f\bigl(\varphi(n-3),\varphi(n),\varphi(n+3)\bigr)\qquad\text{for all }n\in\Z,
\end{equation}
and write $(\mathrm P_n)$ for the instance of \eqref{c79:eq:Phi} at $n$. $\Phi$ is invariant under the shift $n\mapsto n+1$.

\begin{lemma}[coding of $\mathrm{Fix}_1$]\label{c79:lem:phi}
\begin{enumerate}
\item[(a)] The map $\varphi\mapsto c$, $c_i=(\varphi(3i-1),\varphi(3i))$, is a bijection from $\Phi$ onto $\mathrm{Fix}_1$, and $Y_c(t,i)=\varphi(t+3i)$. The sequence of $F^jc$ is $n\mapsto\varphi(n+j)$.
\item[(b)] Let $\varphi,\psi\in\Phi$ code $c,d$. Then $c$ and $d$ agree on a left (right) half-line if and only if $\varphi$ and $\psi$ agree on a left (right) half-line.
\item[(c)] The vacuum $V_\omega$ is coded by $v_\omega(n)=v(n+\omega)$, where $v(n)=[\,n\equiv3,4\pmod 6\,]$; the constant configuration $F^{s}((0,0)^{\Z})$ ($s\in\Z/3$) by $\kappa_s(n)=[\,n\equiv 1-s\pmod 3\,]$; and $(1,1)^{\Z}$ by the constant sequence $1$. The six points of $\mathcal W$ (defined below) lie in $\mathrm{Fix}_1$ and are coded by the shifts of $w(n)=[\,n\not\equiv3,5\pmod 6\,]$.
\end{enumerate}
\end{lemma}

\begin{proof}
(a) $c\in\mathrm{Fix}_1$ means $F^3c=\sh c$, that is $Y_c(t+3,i)=Y_c(t,i+1)$ for all $t,i$, since a configuration is determined by its diagram and $Y_{F^3c}(t,i)=Y_c(t+3,i)$, $Y_{\sh c}(t,i)=Y_c(t,i+1)$. Then $\varphi(n)=Y_c(n,0)$ satisfies $Y_c(t,i)=\varphi(t+3i)$, and \eqref{eq:rel} at $(t,i)$ is $(\mathrm P_{t+3i})$. Conversely, for $\varphi\in\Phi$ the function $Y(t,i)=\varphi(t+3i)$ satisfies \eqref{eq:rel}, so it is the diagram of the configuration $c$ formed by its rows $-1$ and $0$, and $Y(t+3,i)=Y(t,i+1)$ gives $c\in\mathrm{Fix}_1$. The last statement is $Y_{F^jc}(t,i)=Y_c(t+j,i)$.

(b) If $c$ and $d$ agree on $(-\infty,a]$, then, as $F^{\pm1}$ have radius $1$, $Y_c(t,i)=Y_d(t,i)$ for $t\in\{0,1,2\}$ and $i\le a-2$. Since $\varphi(n)=Y_c(n-3i,i)$ for every $i$, taking $i=\lfloor n/3\rfloor$ gives $\varphi(n)=\psi(n)$ for $n\le 3a-4$. Conversely, if $\varphi=\psi$ on $(-\infty,M]$, then $c_i=d_i$ for $3i\le M$. Right half-lines are the same, with $t\in\{-2,-1,0\}$.

(c) By Lemma~\ref{c79:lem:vac}(a), $V(t,i)=V(t+3i,0)=v(t+3i)$, and $F^3V_\omega=\sh V_\omega$; so $V_\omega\in\mathrm{Fix}_1$, coded by $n\mapsto V(n+\omega,0)=v(n+\omega)$. The diagram of $(0,0)^{\Z}$ is $Y(t,i)=[\,t\equiv1\pmod3\,]$, as $F$ acts by $(0,0)^{\Z}\mapsto(0,1)^{\Z}\mapsto(1,0)^{\Z}\mapsto(0,0)^{\Z}$; it is constant in $i$ and of period $3$ in $t$, so it is coded by $\kappa_0$, and $F^s((0,0)^{\Z})$ by $\kappa_s$. The diagram of $(1,1)^{\Z}$ is $1$. For $\mathcal W$: the sequence $w$ is $1,1,1,0,1,0$ on $n\equiv0,\dots,5$, and the configuration it codes has the cells $(w(-1),w(0))=(0,1)$ at even sites and $(w(2),w(3))=(1,0)$ at odd sites. One checks $(\mathrm P_n)$ for $n\equiv0,\dots,5$ with the values of $f$ below; so this point of $\mathcal W$ lies in $\mathrm{Fix}_1$, and its images under $F$ and $\sh$, which are the six points of $\mathcal W$ by the list in the proof of Proposition~\ref{c79:prop:offsets}, are coded by the shifts of $w$.
\end{proof}

We use $f(a,b,0)=1$ and $f(a,b,1)=\neg a\wedge b$ from \eqref{c79:eq:f}. So $(\mathrm P_n)$ says: if $\varphi(n+3)=0$, then $\varphi(n-1)\ne\varphi(n+1)$; if $\varphi(n+3)=1$, then $\varphi(n-1)\oplus\varphi(n+1)=\neg\varphi(n-3)\wedge\varphi(n)$. Read from left to right, this is a rule for $\varphi(n+3)$:

\begin{lemma}\label{c79:lem:Phirules}
Let $\varphi\in\Phi$ and $n,p\in\Z$.
\begin{enumerate}
\item[(a)] \emph{(Forward rule.)} The values $\varphi(n-3),\varphi(n-1),\varphi(n),\varphi(n+1)$ determine $\varphi(n+3)$, unless $\varphi(n-1)\ne\varphi(n+1)$, $\varphi(n-3)=0$ and $\varphi(n)=1$, in which case both values of $\varphi(n+3)$ satisfy $(\mathrm P_n)$. Precisely: if $\varphi(n-1)=\varphi(n+1)$, then $\varphi(n+3)=1$; if $\varphi(n-1)\ne\varphi(n+1)$ and $(\varphi(n-3),\varphi(n))\ne(0,1)$, then $\varphi(n+3)=0$.
\item[(b)] \emph{(Constants absorb.)} If $\varphi$ agrees with $\kappa\in\{\kappa_0,\kappa_1,\kappa_2,1\}$ on a left half-line, then $\varphi=\kappa$.
\item[(c)] \emph{(No new $00$.)} If $\varphi(p)=\varphi(p+1)=0$, then $\varphi(q)=\varphi(q+1)=0$ for some $q\in\{p-1,p-3,p-4\}$. Consequently, if $\varphi(q)\varphi(q+1)\ne00$ for $q=N,\dots,N+3$, then $\varphi(q)\varphi(q+1)\ne00$ for every $q\ge N$.
\item[(d)] If $\varphi=1$ on a right half-line, then $\varphi=1$.
\end{enumerate}
\end{lemma}

\begin{proof}
(a) is the reformulation of $(\mathrm P_n)$ above: if $\varphi(n-1)=\varphi(n+1)$, the value $\varphi(n+3)=0$ is excluded; if $\varphi(n-1)\ne\varphi(n+1)$, the value $\varphi(n+3)=1$ requires $\neg\varphi(n-3)\wedge\varphi(n)=1$, and when this holds both values are allowed, since $f(0,1,c)=1$.

(b) $\kappa$ lies in $\Phi$ by Lemma~\ref{c79:lem:phi}(c), and the exceptional case of (a) never occurs in $\kappa$: in $\kappa_s$ a value $1$ at $n$ has the values $0$ at $n\pm1$, and in $1$ the value at $n-3$ is $1$. So if $\varphi=\kappa$ on $(-\infty,M]$, then by (a) at $n=M-2$ the value $\varphi(M+1)$ is the forced value $\kappa(M+1)$; induct.

(c) By $(\mathrm P_{p-2})$ and $\varphi(p+1)=0$, $\varphi(p-3)\ne\varphi(p-1)$; by $(\mathrm P_{p-3})$ and $\varphi(p)=0$, $\varphi(p-4)\ne\varphi(p-2)$. If $\varphi(p-1)=0$, take $q=p-1$. Otherwise $\varphi(p-1)=1$ and $\varphi(p-3)=0$; if $\varphi(p-2)=0$, take $q=p-3$; otherwise $\varphi(p-2)=1$, so $\varphi(p-4)=0$, and take $q=p-4$. For the consequence, a smallest $p\ge N+4$ with $\varphi(p)\varphi(p+1)=00$ would give such a $q$ with $N\le q<p$.

(d) If $\varphi(p)=0$ and $\varphi(p+2)=\varphi(p+3)=\varphi(p+4)=\varphi(p+6)=1$, then $(\mathrm P_{p+3})$ reads $0=f(0,1,1)=1$. So if $\varphi=1$ on $[M,\infty)$, then $\varphi(M-2)=\varphi(M-1)=1$ (take $p=M-2$ and $p=M-1$); induct.
\end{proof}

Let $\varphi_K\in\{0,1\}^{\Z}$ be given by
\[
\varphi_K(n)=\begin{cases}v(n)=[\,n\equiv3,4\pmod 6\,]&(n\le6),\\ \kappa_0(n)=[\,n\equiv1\pmod 3\,]&(n\ge4);\end{cases}
\]
the two formulas agree on $\{4,5,6\}$, and $\varphi_K(n)$ for $n=1,\dots,10$ is $0,0,1,1,0,0,1,0,0,1$. Let $K^-$ be the configuration equal to $V_0$ on $(-\infty,1]$ and to $(0,0)$ on $[2,\infty)$, and $K^+$ the configuration equal to $V_0$ on $(-\infty,0]$ and to $(0,1)$ on $[1,\infty)$. For $j\in\Z$ put $K^-_j=F^jK^-$ and $K^+_j=F^jK^+$. Let $\mathcal W$ be the set of the six configurations of period $2$ whose cells at an even and at an odd site are two distinct elements of $\{(0,1),(1,0),(1,1)\}$.

\begin{lemma}[kinks]\label{c79:lem:kinks}
\begin{enumerate}
\item[(a)] $\varphi_K\in\Phi$, and $\varphi_K$ codes $K^-$; so $Y_{K^-}(t,i)=\varphi_K(t+3i)$. Moreover $K^+=F^{-2}\rho K^-$, so $Y_{K^+}(t,i)=\varphi_K(1-t+3i)$. Hence $F^3K^-=\sh K^-$, $F^3K^+=\sh^{-1}K^+$ and $\rho K^-=F^2K^+$; so $K^-_j\in\mathrm{Fix}_1$, $K^+_j\in\mathrm{Fix}_{-1}$, $\sh K^-_j=K^-_{j+3}$, $\sh^{-1}K^+_j=K^+_{j+3}$, $\rho K^-_j=K^+_{2-j}$ and $\rho K^+_j=K^-_{2-j}$.
\item[(b)] $K^-_j$ agrees with $V_j$ exactly on $(-\infty,p^-(j))$, and $K^+_j$ with $V_j$ exactly on $(-\infty,q^+(j))$, where $p^-(j)=2-\lfloor(j-1)/3\rfloor$ and $q^+(j)=2+\lfloor(j+1)/3\rfloor$. The configuration $K^-_j$ is constant, equal to $F^j((0,0)^{\Z})$, exactly on $[2-\lfloor(j+1)/3\rfloor,\infty)$, and $K^+_j$ is constant, equal to $F^j((0,1)^{\Z})$, on the cells $i$ with $3i\ge j+3$.
\item[(c)] The points of $\mathrm{Fix}_1$ that agree with a vacuum on a left half-line and with a constant configuration on a right half-line are the $K^-_j$, $j\in\Z$, and these are pairwise distinct; the same holds for $\mathrm{Fix}_{-1}$ and the $K^+_j$. No point of $\mathrm{Fix}_{\pm1}$ agrees with a point of $\mathcal W$ on a left half-line and with a constant configuration on a right half-line, and a point of $\mathrm{Fix}_{\pm1}$ that agrees with constant configurations on a left and on a right half-line is constant.
\end{enumerate}
\end{lemma}

\begin{proof}
(a) $(\mathrm P_n)$ holds for $n\le3$, where it involves only values of $v\in\Phi$, and for $n\ge7$, where it involves only values of $\kappa_0\in\Phi$. The remaining three instances are $(\mathrm P_4)$: $\varphi_K(3)\oplus\varphi_K(5)=1=f(0,1,1)$; $(\mathrm P_5)$: $\varphi_K(4)\oplus\varphi_K(6)=1=f(0,0,0)$; and $(\mathrm P_6)$: $\varphi_K(5)\oplus\varphi_K(7)=1=f(1,0,0)$. The configuration coded by $\varphi_K$ has the cells $(\varphi_K(3i-1),\varphi_K(3i))$, which are $(v(3i-1),v(3i))=(V_0)_i$ for $i\le1$ (Lemma~\ref{c79:lem:phi}(c)) and $(\kappa_0(3i-1),\kappa_0(3i))=(0,0)$ for $i\ge2$; this is $K^-$. By \eqref{eq:rho}, the diagram of $F^{-2}\rho K^-$ is $Y_{K^-}(1-t,i)=\varphi_K(1-t+3i)$, and its cells are $(\varphi_K(2+3i),\varphi_K(1+3i))$. For $i\le0$ these are $(v(2+3i),v(1+3i))=(v(3i-1),v(3i))=(V_0)_i$, by $v(1-n)=v(n)$ and the period $6$; for $i\ge1$ they are $(\kappa_0(2+3i),\kappa_0(1+3i))=(0,1)$. So $F^{-2}\rho K^-=K^+$. Now $K^-\in\mathrm{Fix}_1$, and $F^3K^+=F^{-2}\rho F^{-3}K^-=F^{-2}\rho\sh^{-1}K^-=\sh^{-1}K^+$; the rest follows, as $\sh$ commutes with $F$ and $\rho$, and $\rho F^j=F^{-j}\rho$.

(b) By (a), $(K^-_j)_i=(\varphi_K(j-1+3i),\varphi_K(j+3i))$ and $(V_j)_i=(v(j-1+3i),v(j+3i))$. As $\varphi_K=v$ on $(-\infty,6]$ while $\varphi_K(7)=1\ne v(7)$ and $\varphi_K(9)=0\ne v(9)$, the two cells agree if $j+3i\le6$ and differ if $j+3i\in\{7,8,9\}$ (one of the two arguments is then $7$ or $9$). So $p^-(j)$ is the least $i$ with $j+3i\ge7$, which is $\lceil(7-j)/3\rceil=2-\lfloor(j-1)/3\rfloor$. In the same way $(K^+_j)_i=(\varphi_K(2-j+3i),\varphi_K(1-j+3i))$, the vacuum cell is the same with $v$, and $q^+(j)$ is the least $i$ with $2-j+3i\ge7$, which is $2+\lfloor(j+1)/3\rfloor$. If $j-1+3i\ge4$, then $(K^-_j)_i=(\kappa_0(j-1),\kappa_0(j))$, the cell of $F^j((0,0)^{\Z})$ by Lemma~\ref{c79:lem:phi}(c); and if $1-j+3i\ge4$, then $(K^+_j)_i=(\kappa_0(2-j),\kappa_0(1-j))=(\kappa_0(j),\kappa_0(j+1))$, the cell of $F^{j+1}((0,0)^{\Z})=F^j((0,1)^{\Z})$. The least $i$ with $j-1+3i\ge4$ is $i_0=2-\lfloor(j+1)/3\rfloor$, and $j-1+3(i_0-1)=e+1$ with $e\in\{0,1,2\}$, $e\equiv j+1\pmod 3$; the cell $(K^-_j)_{i_0-1}=(\varphi_K(e+1),\varphi_K(e+2))$ is $(0,0)$, $(0,1)$, $(1,1)$ for $e=0,1,2$, while the constant cell $(\kappa_0(e+1),\kappa_0(e+2))$ is $(1,0)$, $(0,0)$, $(0,1)$. So $K^-_j$ is not constant on $[i_0-1,\infty)$.

(c) By Lemma~\ref{c79:lem:phi}, a point $c\in\mathrm{Fix}_1$ that agrees with a vacuum on a left half-line and with a constant configuration on a right half-line is coded by a $\varphi\in\Phi$ that agrees with some $v_\omega$ on a left half-line and with some $\kappa\in\{\kappa_0,\kappa_1,\kappa_2,1\}$ on a right half-line. By Lemma~\ref{c79:lem:Phirules}(d), $\kappa\ne1$, since $\varphi\ne1$. So $\varphi\ne v_\omega$, and, replacing $\varphi$ by a shift, we may assume that $\omega=0$ and that $N$ is the least integer with $\varphi(N)\ne v(N)$. The values $\varphi(N-6),\varphi(N-4),\varphi(N-3),\varphi(N-2)$ are those of $v$, and $v(N)$ satisfies $(\mathrm P_{N-3})$; by Lemma~\ref{c79:lem:Phirules}(a), $\varphi(N)\ne v(N)$ is possible only if $v(N-3)=1$, $v(N-6)=0$ and $v(N-4)\ne v(N-2)$. This means $N\equiv0$ or $1\pmod 6$ (then $v(N-6)=v(N)=0$ and $v(N-4)\ne v(N-2)$ hold), and $\varphi(N)=1$.

If $N\equiv0$, then $\varphi(N-5),\dots,\varphi(N)$ are $0,0,1,1,0,1$, and by Lemma~\ref{c79:lem:Phirules}(a) at $n=N-1$, where $\varphi(N-2)=\varphi(N)=1$, $\varphi(N+2)=1$. So no $00$ starts at $N-2,N-1,N,N+1$, and by Lemma~\ref{c79:lem:Phirules}(c) none starts at any $q\ge N-2$; but $\kappa_s$ contains $00$ in every interval of length $3$, a contradiction. Hence $N\equiv1$, and $\varphi(N-6),\dots,\varphi(N)$ are $0,0,1,1,0,0,1$. Lemma~\ref{c79:lem:Phirules}(a) at $n=N-2$ (values $0,1,0,0$) gives $\varphi(N+1)=0$, at $n=N-1$ (values $1,0,0,1$) gives $\varphi(N+2)=0$, and at $n=N$ (values $1,0,1,0$) gives $\varphi(N+3)=1$. So $\varphi=\kappa$ on $[N-2,N+3]$ for the $\kappa\in\{\kappa_0,\kappa_1,\kappa_2\}$ with $\kappa(N)=1$, and by the argument of Lemma~\ref{c79:lem:Phirules}(b), applied from $N+3$ on, $\varphi=\kappa$ on $[N-2,\infty)$. Thus $\varphi=v$ on $(-\infty,N-1]$ and $\varphi=\kappa$ on $[N-2,\infty)$, with $N\equiv1\pmod6$; this is the shift $n\mapsto\varphi_K(n+7-N)$ of $\varphi_K$, which codes $F^{7-N}K^-$ by Lemma~\ref{c79:lem:phi}(a). Undoing the first shift, $c$ is some $K^-_j$. The $K^-_j$ are pairwise distinct, since they are coded by the shifts of $\varphi_K$, which is not periodic.

If $c\in\mathrm{Fix}_1$ agrees with a point of $\mathcal W$ on a left half-line, then $\varphi$ agrees with a shift of $w$ on a left half-line (Lemma~\ref{c79:lem:phi}(b),(c)). As $w$ contains no $00$, Lemma~\ref{c79:lem:Phirules}(c) shows that $\varphi$ contains no $00$, so it does not agree with any $\kappa_s$ on a right half-line, and by Lemma~\ref{c79:lem:Phirules}(d) it does not agree with $1$ on a right half-line, since $\varphi\ne1$. If $c$ agrees with a constant configuration on a left half-line, then $\varphi$ agrees with some $\kappa$ on a left half-line, and $\varphi=\kappa$ by Lemma~\ref{c79:lem:Phirules}(b), so $c$ is constant.

For $\mathrm{Fix}_{-1}$: since $\rho\sh^{-1}F^3\rho=\sh^{-1}F^{-3}$, whose fixed points are those of $\sh F^3$, $\rho$ maps $\mathrm{Fix}_1$ bijectively onto $\mathrm{Fix}_{-1}$. It acts cellwise, maps the vacua onto the vacua ($\rho V_\omega=V_{2-\omega}$), the constant configurations onto the constant configurations, $\mathcal W$ onto $\mathcal W$, and $K^-_j$ to $K^+_{2-j}$ by (a). So the statements for $\mathrm{Fix}_{-1}$ are those for $\mathrm{Fix}_1$.
\end{proof}

\begin{proposition}[normalisation]\label{c79:prop:offsets}
Let $g\in\Aut(X_F)$. There are $j\in\{0,1,2\}$ and $a,b\in\Z$ such that $h=F^{-3b}\sh^{-a}F^{-j}g$ fixes $K^-$ and $K^+$.
\end{proposition}

\begin{proof}
The automorphism $g$ permutes the fixed points of $\sh$, which are the four constant configurations, and the induced permutation commutes with that of $F$, the $3$-cycle $(0,0)^{\Z}\mapsto(0,1)^{\Z}\mapsto(1,0)^{\Z}$. The centralizer of a $3$-cycle in $\Sym(4)$ is the group it generates, so $g$ acts on the constants as $F^j$ for one $j\in\{0,1,2\}$, and $g'=F^{-j}g$ fixes every constant configuration. It permutes the fixed points of $\sh^2$, which are the configurations of period $2$, and preserves the sizes of their $F$-orbits. By Lemma~\ref{c79:lem:vac}(a) the vacua form one $F$-orbit of size $6$, and the only other orbit of size $6$ is $\mathcal W$: since $f(p,q,p)=\neg p$, $F$ maps the configuration of period $2$ with cells $(x_e,y_e)$ at even and $(x_o,y_o)$ at odd sites to the one with cells $(y_e,x_e\oplus\neg y_o)$ and $(y_o,x_o\oplus\neg y_e)$, and its orbits on the sixteen configurations of period $2$, written as pairs (even cell, odd cell) in the coding $2x+y$, are $\{(0,0),(1,1),(2,2)\}$, $\{(3,3)\}$, the vacua $(0,1)\mapsto(0,3)\mapsto(0,2)\mapsto(1,0)\mapsto(3,0)\mapsto(2,0)\mapsto(0,1)$, and $\mathcal W$: $(1,2)\mapsto(3,1)\mapsto(3,2)\mapsto(2,1)\mapsto(1,3)\mapsto(2,3)\mapsto(1,2)$. Since $K^-$ agrees with $V_0$ on a left half-line and with $(0,0)^{\Z}$ on a right half-line, and $g'$ is a block map, $g'K^-\in\mathrm{Fix}_1$ agrees with $g'V_0$ on a left half-line and with $(0,0)^{\Z}$ on a right half-line. By Lemma~\ref{c79:lem:kinks}(c), $g'V_0\notin \mathcal W$, so $g'V_0=V_{\omega_0}$ for some $\omega_0$, and $g'K^-=K^-_{j'}$ for some $j'$. Its right constant is $F^{j'}((0,0)^{\Z})=(0,0)^{\Z}$, so $j'=3e^-$ for an integer $e^-$, and $g'K^-=\sh^{e^-}K^-$. In the same way, $g'K^+\in\mathrm{Fix}_{-1}$ agrees with $g'V_0$ on a left half-line and with $(0,1)^{\Z}$ on a right half-line; by the statements of Lemma~\ref{c79:lem:kinks}(c) for $\mathrm{Fix}_{-1}$, $g'V_0\notin\mathcal W$ and $g'K^+=K^+_{j''}$ for some $j''$, whose right constant $F^{j''}((0,1)^{\Z})=(0,1)^{\Z}$ gives $j''=3n$ for an integer $n$, so $g'K^+=K^+_{3n}=\sh^{-n}K^+$; put $e^+=-n$, so $g'K^+=\sh^{e^+}K^+$. Moreover, $g'K^-$ agrees on a left half-line with $g'V_0=V_{\omega_0}$ and also with $\sh^{e^-}V_0=V_{3e^-}$; so $\omega_0\equiv3e^-$, and likewise $\omega_0\equiv3e^+\pmod 6$. Hence $e^-\equiv e^+\pmod2$. Put $a=\frac12(e^-+e^+)$ and $b=\frac12(e^--e^+)$. Since $F^3$ fixes the constants and $F^3K^\mp=\sh^{\pm1}K^\mp$, the map $\sh^aF^{3b}$ sends $K^-$ to $\sh^{a+b}K^-=\sh^{e^-}K^-$ and $K^+$ to $\sh^{a-b}K^+=\sh^{e^+}K^+$, so $h=F^{-3b}\sh^{-a}g'$ fixes $K^-$ and $K^+$.
\end{proof}

\begin{lemma}[Lemma K$_{93}$]\label{c79:lem:K}
Let $h$ be a block map that commutes with $F_{93}$ and fixes $K^-$ and $K^+$. Then $h=\mathrm{id}$.
\end{lemma}

Lemma~K$_{93}$ is proved in Section~\ref{c79:sec:end}. Let $\mathcal H$ be the set of block maps $h$, of any radius, that commute with $F$ and fix $K^-$ and $K^+$. Such an $h$ fixes every $K^\pm_j$. It fixes every vacuum: $hV_0$ has period $2$ and agrees with $V_0$ on a left half-line, because $hK^-=K^-$; so $hV_0=V_0$, and $hV_\omega=V_\omega$ by $F$. It fixes $(0,0)^{\Z}$, $(0,1)^{\Z}$ and $(1,0)^{\Z}$, the right constants of the $K^\pm_j$. And $\rho h\rho\in\mathcal H$, since $\rho F\rho=F^{-1}$ and $\rho h\rho K^\mp=\rho hF^2K^\pm=\rho F^2K^\pm=K^\mp$ by Lemma~\ref{c79:lem:kinks}(a).

\begin{lemma}\label{c79:lem:hglue}
Let $h$ be a block map of radius $R$ that fixes $V_\omega$ and $\Lambda$, let $c$ agree with $V_\omega$ on $(r,\infty)$ and $\Lambda$ with $V_\omega$ on $(-\infty,p)$, where $p\ge r+2R+1$. Then $hc$ agrees with $V_\omega$ on $(r+R,\infty)$ and $h(c\sqcup_r\Lambda)=hc\sqcup_{r+R}\Lambda$.
\end{lemma}

\begin{proof}
For $i\ge r+R+1$ the window of $i$ lies in $(r,\infty)$, where $c=V_\omega$, so $(hc)_i=(hV_\omega)_i=(V_\omega)_i$. Put $\hat c=c\sqcup_r\Lambda$. For $i\le p-1-R$ the window of $i$ lies in $(-\infty,p-1]$, where $\hat c=c$, so $(h\hat c)_i=(hc)_i$; for $i\ge r+R+1$ it lies in $(r,\infty)$, where $\hat c=\Lambda$, so $(h\hat c)_i=(h\Lambda)_i=\Lambda_i$. These ranges cover $\Z$, and on their overlap $hc$ and $\Lambda$ both agree with $V_\omega$.
\end{proof}

\subsection{Moves}\label{c79:sec:moves}

\begin{definition}\label{c79:def:moves}
Let $c\in\Omega_r$ have phase $\omega$, let $m\in\Z$ and $\omega'\in\Z/6$. Let $S^{\omega'}_m(c)$ be the configuration $D$ with $F^mD=F^mc\sqcup_rV_{\omega'+m}$, that is, $F^mD$ agrees with $F^mc$ on $(-\infty,r]$ and with the vacuum of phase $\omega'+m$ on $(r,\infty)$. Put $\vartheta=\vartheta(c)$ and $z_t=Y_c(t,r)$, and define the \emph{move at time $m$}
\[
T_m(x)=\begin{cases}S_m^{\omega+2}(x)&\text{if }m+\vartheta\equiv2\text{ and }z_{m-1}=0,\\ S_m^{\omega-1}(x)&\text{if }m+\vartheta\equiv2\text{ and }z_{m-1}=1,\\ S_m^{\omega-2}(x)&\text{if }m+\vartheta\equiv0\text{ and }z_{m}=0,\\ S_m^{\omega+1}(x)&\text{if }m+\vartheta\equiv0\text{ and }z_{m}=1,\\ x&\text{otherwise,}\end{cases}
\]
congruences modulo $6$. We call $T_m$ a \emph{forward move} in the first two cases and a \emph{backward move} in the next two, and write $M_m$, $\bar M_m$ for $T_m$ restricted to the configurations on which it is forward, respectively backward.
\end{definition}

So a move keeps the state at time $m$ on the columns $\le r$ and switches the phase of the vacuum to its right, by $+2$ or $-1$ (forward) or by $-2$ or $+1$ (backward), according to one bit of the edge column.

\begin{lemma}\label{c79:lem:Tbasic}
Let $c\in\Omega_r$ and $m\in\Z$.
\begin{enumerate}
\item[(a)] $T_mc\in\Omega_r$.
\item[(b)] $T_mc$ agrees with $c$ on $(-\infty,r-|m|]$; and $Y_{T_mc}(t,r)=z_t$ for $t\in\{m-1,m\}$.
\item[(c)] $T_m(c)=F^{-\ell}T_{m-\ell}(F^\ell c)$ for every $\ell\in\Z$, and $\sh T_m=T_m\sh$ (with edges $r$ and $r-1$). Moreover $\rho c\in\Omega_r$, $\vartheta(\rho c)=2-\vartheta(c)$, the edge column of $\rho c$ is $t\mapsto z_{-1-t}$, and $\bar M_m=\rho M_{-m}\rho$. So under $\rho$ the time $t$ of a diagram corresponds to the time $-1-t$, while the move at time $m$ corresponds to the move at time $-m$.
\end{enumerate}
\end{lemma}

\begin{proof}
(a) By Lemma~\ref{c79:lem:glue}(b) it suffices that the cell $r$ of $F^mc$, which is $(z_{m-1},z_m)$, differs from the cell of $V_{\omega'+m}$ at $r$, which by Lemma~\ref{c79:lem:glue}(c) has bits $[m-1+\vartheta'\equiv3,4]$ and $[m+\vartheta'\equiv3,4]$, where $\vartheta'=\omega'+3r$. In the four cases of the definition, $m+\vartheta'\equiv4,1,4,1$, so this cell is $(1,1)$, $(0,0)$, $(1,1)$, $(0,0)$, while $(z_{m-1},z_m)$ is $(0,\cdot)$, $(1,\cdot)$, $(\cdot,0)$, $(\cdot,1)$.

(b) $F^mT_mc$ and $F^mc$ agree on $(-\infty,r]$, and $F^{-m}$ has radius $|m|$. The second statement holds because rows $m-1$ and $m$ form the state at time $m$.

(c) $F^\ell c\in\Omega_r$ has phase $\omega+\ell$, so $\vartheta(F^\ell c)=\vartheta+\ell$, its edge column is $t\mapsto z_{t+\ell}$, and the case of $T_{m-\ell}$ at $F^\ell c$ is the case of $T_m$ at $c$; and $F^{m-\ell}(F^\ell S^{\omega'}_mc)=F^mc\sqcup_rV_{\omega'+m}$, which says $F^\ell S^{\omega'}_m(c)=S^{\omega'+\ell}_{m-\ell}(F^\ell c)$. The shift: $\sh c\in\Omega_{r-1}$ has phase $\omega+3$ and the same $\vartheta$ and edge column. For $\rho$: by Lemma~\ref{c79:lem:vac}(a) and \eqref{eq:rho}, $\rho c\in\Omega_r$ has phase $2-\omega$, $\vartheta(\rho c)=2-\vartheta$, and edge column $t\mapsto z_{-1-t}$; and $\rho S^{\omega'}_{m'}\rho=S^{2-\omega'}_{-m'}$, since $F^{-m'}\rho=\rho F^{m'}$ exchanges the two rows of each cell. The forward move $M_{-m}$ at $\rho c$ applies when $-m+2-\vartheta\equiv2$, that is $m+\vartheta\equiv0$, reads $Y_{\rho c}(-m-1,r)=z_m$, and gives the phase $2-\omega+2$ or $2-\omega-1$; conjugating by $\rho$ gives $S_m^{\omega-2}$ or $S_m^{\omega+1}$, which is $\bar M_m$.
\end{proof}

A move is not undone by a move at the same time: after $T_m$ the class of $m+\vartheta$ is $4$ or $1$, and $T_m$ acts as the identity. Orbits of moves are therefore defined below by directed reachability.

The moves are realised by gluing a kink. If $c\in\Omega_r$ has phase $\omega$ and $j\equiv\omega\pmod 6$, $j\le-3r$, then $K^-_j$ agrees with $V_j=V_\omega$ on $(-\infty,r+2]$ by Lemma~\ref{c79:lem:kinks}(b), and $c\sqcup_rK^-_j$ is defined. We say that a configuration $\hat c$ has \emph{outgoing index} $J\in\Z$ if there are $q\in\Z$ and $s_0$ such that $F^s{\hat c}$ agrees with $K^+_{s+J}$ on $[q,\infty)$ for every $s\ge s_0$. The index is unique if it exists: for $J\not\equiv J'\pmod 3$ the right constants of $K^+_{s+J}$ and $K^+_{s+J'}$ differ; for $J'=J+3d$ with $d$ odd, $K^+_{s+J'}=\sh^{-d}K^+_{s+J}$ and the vacua $V_{s+J}$, $V_{s+J'}$ to the left of the two kinks differ at every site; and for $d\ne0$ even the two vacua coincide while the first sites $q^+(s+J)\ne q^+(s+J')$ at which the configurations leave them both tend to $\infty$. For the same reason, $K^-_j$ has no outgoing index: $F^sK^-_j=K^-_{j+s}$ is constant on $[q,\infty)$ for large $s$, while $K^+_{s+J}$ is not, as no two adjacent cells of a vacuum are equal.

\begin{lemma}[moves from kinks]\label{c79:lem:probe}
Let $c\in\Omega_r$ have phase $\omega$, let $j\equiv\omega\pmod 6$ with $j\le-3r$, and put $\hat c=c\sqcup_rK^-_j$, $m=2-j-3r$, $z=Y_c(m-1,r)$, $D=M_mc$ and
\[
J=j+6r+3z-4.
\]
Then $m\ge2$, $m+\vartheta(c)\equiv2\pmod6$, and for every $s\ge m$ the configuration $F^s\hat c$ agrees with $F^sD$ on $(-\infty,r]$ and with $K^+_{s+J}$ on $(r,\infty)$:
\[
F^s\hat c=F^sD\sqcup_rK^+_{s+J}\qquad(s\ge m).
\]
In particular $\hat c$ has outgoing index $J$.
\end{lemma}

\begin{proof}
$m+\vartheta\equiv2-j-3r+\omega+3r\equiv2$, and $m\ge2$ as $j\le-3r$. The phase of $D$ is $\omega'=\omega+2-3z\equiv J\pmod 6$. We write down the diagram of $\hat c$. Put $\tau=t-m$, $k=i-r$, and let
\begin{align*}
A(t,i)&=Y_{K^-_j}(t,i)=\varphi_K(t+j+3i)=\varphi_K(\tau+3k+2),\\
B(t,i)&=Y_{K^+_J}(t,i)=\varphi_K(1-t-J+3i)=\varphi_K(3k-\tau+3-3z),
\end{align*}
by Lemma~\ref{c79:lem:kinks}(a). Both are functions of $(t,i)$; written in $\tau$ and $k$, neither depends on $r$, $\omega$ or $j$, and below we say that a statement holds at $(\tau,k)$ when it holds at $(t,i)=(m+\tau,r+k)$. Define $Z\colon\Z^2\to\{0,1\}$ by
\[
Z(t,i)=\begin{cases}
Y_c(t,i)&i\le r,\ t\le m,\\
Y_D(t,i)&i\le r,\ t\ge m-1,\\
A(t,i)&i>r,\ t\le m-1,\\
B(t,i)&i>r,\ t\ge m.
\end{cases}
\]
The first two lines are consistent, since $F^mD$ agrees with $F^mc$ on $(-\infty,r]$, and the rows $m-1$, $m$ form the state at time $m$. Recall $\varphi_K(n)=v(n)$ for $n\le6$, $\varphi_K(n)=\kappa_0(n)=[n\equiv1\ (3)]$ for $n\ge4$, and $\varphi_K(1),\dots,\varphi_K(9)=0,0,1,1,0,0,1,0,0$. We use four facts.
\begin{enumerate}
\item[(i)] $A(t,i)=V(t+\omega,i)$ whenever $\tau+3k+2\le6$, in particular for $k=1$, $\tau\le1$ and for $k=2$, $\tau\le-2$; and $B(t,i)=V(t+J,i)=V(t+\omega',i)$ whenever $3k-\tau+3-3z\le6$, in particular for $k=1$, $\tau\ge-3z$ and for $k=2$, $\tau\ge3-3z$. (By Lemma~\ref{c79:lem:phi}(c), $V(t,i)=v(t+3i)$, and $v(1-n)=v(n)$; and $j\equiv\omega$, $J\equiv\omega'$.)
\item[(ii)] $A(t,i)=B(t,i)$ for $\tau=-1$, $k\ge1$, and for $\tau=0$, $k\ge2$. (For $\tau=-1$ the arguments $3k+1$ and $3k+4-3z$ are $\ge4$ and $\equiv1\pmod3$, so both values are $1$; for $\tau=0$ the arguments $3k+2\ge8$ and $3k+3-3z\ge6$ are $\not\equiv1\pmod3$, so both values are $0$.)
\item[(iii)] For $k=1$ and $\tau\ge0$, $B(t,r+1)=B(t,r+2)=1$ only if $(\tau,z)=(2,0)$; for $k=1$ and $\tau\le-2$, $A(t,r+1)=A(t,r+2)=1$ never holds. (In the ranges of (i) these are two values of a row of $V$ at adjacent sites, which are never both $1$ by Lemma~\ref{c79:lem:vac}(c); the remaining values are $B(m,r+1)=\varphi_K(6-3z)$ and $B(m,r+2)=\varphi_K(9-3z)=0$, $B(m+1,r+1)=\varphi_K(5-3z)=0$, and, for $z=0$, $B(m+2,r+1)=\varphi_K(4)=1$ and $B(m+2,r+2)=\varphi_K(7)=1$.)
\item[(iv)] If $z=0$, then $Y_D(m+2,r)=0$. (The cell $(z_{m-1},z_m)$ of $F^mc$ at $r$ differs from the cell of $V_{\omega+m}$ at $r$, which by Lemma~\ref{c79:lem:glue}(c) is $(0,0)$, as $m+\vartheta\equiv2$; so $z_m=1$. $D$ has $\vartheta(D)=\vartheta+2$, so $(m+1)+\vartheta(D)\equiv5$, and Lemma~\ref{c79:lem:glue}(c) for $D$ gives $Y_D(m+2,r)=\neg Y_D(m,r)=\neg z_m=0$.)
\end{enumerate}
We check \eqref{eq:rel} for $Z$ at every $(t,i)$; it involves $Z$ at $(t\pm1,i)$ and $(t,i-1),(t,i),(t,i+1)$.

\emph{Columns $i\le r-1$.} For $t\le m-1$ all values are those of $Y_c$, for $t\ge m$ all are those of $Y_D$.

\emph{Column $r$.} For $t\le m-1$ the values on the columns $\le r$ are those of $Y_c$, and $Z(t,r+1)=A(t,r+1)=V(t+\omega,r+1)$ by (i), since $\tau\le-1$; this is \eqref{eq:rel} for $Y_c$, whose column $r+1$ is $V(t+\omega,r+1)$ (Lemma~\ref{c79:lem:glue}(b)). For $t\ge m$ they are those of $Y_D$, and $Z(t,r+1)=B(t,r+1)=V(t+\omega',r+1)$ by (i), since $\tau\ge0\ge-3z$; this is \eqref{eq:rel} for $Y_D$.

\emph{Columns $i\ge r+2$.} For $t\le m-2$ all values are values of $A$, and for $t\ge m+1$ all are values of $B$. For $t=m-1$ they are values of $A$, except $Z(m,i)=B(m,i)$, which equals $A(m,i)$ by (ii); for $t=m$ they are values of $B$, except $Z(m-1,i)=A(m-1,i)=B(m-1,i)$ by (ii). In each case \eqref{eq:rel} is that of the diagram $A$ or $B$.

\emph{Column $r+1$.} Here $f(Z(t,r),Z(t,r+1),Z(t,r+2))$ depends on $Z(t,r)$ only if $Z(t,r+1)=Z(t,r+2)=1$. For $t\le m-2$ all other values are values of $A$, and by (iii) the relation is that of $A$. For $t\ge m$ all other values are values of $B$ (at $t=m$, $Z(m-1,r+1)=B(m-1,r+1)$ by (ii)), and by (iii) the relation is that of $B$, except at $(\tau,z)=(2,0)$. There it reads $B(m+3,r+1)\oplus B(m+1,r+1)=\varphi_K(3)\oplus\varphi_K(5)=1=f(Z(m+2,r),1,1)=\neg Y_D(m+2,r)$, which is (iv). At $t=m-1$ the values are $Z(m-1,r)=z$, $Z(m-1,r+1)=A(m-1,r+1)=\varphi_K(4)=1$, $Z(m-1,r+2)=\varphi_K(7)=1$, $Z(m-2,r+1)=\varphi_K(3)=1$ and $Z(m,r+1)=B(m,r+1)=\varphi_K(6-3z)=z$, and the relation reads $z\oplus1=f(z,1,1)=\neg z$.

So $Z$ satisfies \eqref{eq:rel} everywhere. Its rows $-1$ and $0$ are those of $c$ on $(-\infty,r]$ and of $K^-_j$ on $(r,\infty)$, as $m\ge2$; so $Z=Y_{\hat c}$. For $s\ge m$, rows $s-1$ and $s$ of $Z$ are those of $Y_D$ on the columns $\le r$, and those of $B$ on the columns $>r$ (for $s=m$, by (ii) at $\tau=-1$); this is the claim. Finally $\hat c$ has outgoing index $J$, with $q=r+1$ and $s_0=m$.
\end{proof}

The lemma holds, with the same proof, for every $j\equiv\omega$ with $j\le1-3r$ (then $m\ge1$). In the coordinates $(\tau,k)$, the incoming kink leaves the vacuum of phase $\omega$ on the line $\tau+3k+2=7$, the outgoing kink leaves the vacuum of phase $\omega'$ on the line $3k-\tau+3-3z=7$, and the edge bit $z=Y_c(m-1,r)$ is read at the single point $(\tau,k)=(-1,1)$, through $f(z,1,1)=\neg z$.

By Lemma~\ref{c79:lem:Tbasic}(c), the backward moves are realised in the same way by $\rho$: $\bar M_m=\rho M_{-m}\rho$, and gluing $K^-_j$ to $\rho c$ is gluing $\rho K^-_j=K^+_{2-j}$ to $c$ in reversed time.

\subsection{The first columns}\label{c79:sec:H}

For $k\ge1$ consider
\begin{itemize}
\item[$(\mathrm H_k)$] \emph{for every $h\in\mathcal H$, every $r$ and every $c\in\Omega_r$, $hc\in\Omega_r$ has the phase of $c$, and $Y_{hc}(t,i)=Y_c(t,i)$ for all $t\in\Z$ and $i\ge r-k+1$.}
\end{itemize}

\begin{proposition}[$(\mathrm H_1)$]\label{c79:prop:H1}
$(\mathrm H_1)$ holds.
\end{proposition}

\begin{proof}
Let $h\in\mathcal H$ have radius $R$, let $c\in\Omega_r$ have phase $\omega$, and write $z_t=Y_c(t,r)$, $\vartheta=\vartheta(c)$. By Lemma~\ref{c79:lem:hglue} with $\Lambda=V_\omega$, $hc$ agrees with $V_\omega$ on $(r+R,\infty)$. Let $j\equiv\omega\pmod 6$ with $j\le-3(r+2R)$, so that $p^-(j)\ge r+2R+3$, and $\hat c=c\sqcup_rK^-_j$. By Lemma~\ref{c79:lem:hglue}, $h{\hat c}=hc\sqcup_{r+R}K^-_j$. If $hc=V_\omega$, then $h{\hat c}=K^-_j$, which has no outgoing index; but $F^sh{\hat c}=hF^s{\hat c}$ agrees with $hK^+_{s+J}=K^+_{s+J}$ on $[q+R,\infty)$ whenever $F^s{\hat c}$ agrees with $K^+_{s+J}$ on $[q,\infty)$, so $h{\hat c}$ has the outgoing index $J$ of $\hat c$, which exists by Lemma~\ref{c79:lem:probe}. Hence $hc\in\Omega_{r'}$ with phase $\omega$, for some $r'\le r+R$, and $h{\hat c}=hc\sqcup_{r'}K^-_j$ with $j\le-3r'$. Lemma~\ref{c79:lem:probe} for $c$ and for $hc$, with $m=2-j-3r$ and $m'=2-j-3r'$, gives for the common outgoing index $J$ of $\hat c$ and $h\hat c$
\[
j+6r'+3Y_{hc}(m'-1,r')-4=J=j+6r+3z_{m-1}-4.
\]
So $6(r'-r)\in\{-3,0,3\}$, hence $r'=r$, $m'=m$ and $Y_{hc}(m-1,r)=z_{m-1}$. As $j$ ranges over $j\equiv\omega$, $j\le-3(r+2R)$, the time $t=m-1$ ranges over all $t\equiv1-\vartheta\pmod6$ with $t\ge t_0$, for some $t_0$. Applying this to $F^{-6i}c$, which lies in $\Omega_r$ with phase $\omega$ and satisfies $hF^{-6i}c=F^{-6i}hc$, gives $Y_{hc}(t,r)=z_t$ for every $t\equiv1-\vartheta$. Next apply the same to $\rho h\rho\in\mathcal H$ and $\rho c\in\Omega_r$, which has $\vartheta(\rho c)=2-\vartheta$ and edge column $t\mapsto z_{-1-t}$ (Lemma~\ref{c79:lem:Tbasic}(c)): $Y_{\rho hc}(t,r)=Y_{\rho c}(t,r)$ for $t\equiv\vartheta-1$, that is $Y_{hc}(t,r)=z_t$ for $t\equiv-\vartheta$. Both $c$ and $hc$ lie in $\Omega_r$ with phase $\omega$, so by Lemma~\ref{c79:lem:glue}(c) both edge columns satisfy $z_{t+1}=\neg z_{t-1}$ whenever $t+\vartheta\equiv2,3,4,5$. Writing $t+\vartheta\equiv n$, the values with $n=0,1$ are known; $n=3$ follows from $n=1$, $n=4$ from $n=0$ (at $t+\vartheta\equiv5$), $n=2$ from $n=4$, and $n=5$ from $n=3$. So the edge columns agree at all times, and to the right of $r$ both diagrams are the vacuum of phase $\omega$.
\end{proof}

\begin{lemma}[moves commute with $\mathcal H$]\label{c79:lem:commute}
For every $h\in\mathcal H$, $r\in\Z$, $c\in\Omega_r$ and $m\in\Z$, $hT_mc=T_mhc$.
\end{lemma}

\begin{proof}
By $(\mathrm H_1)$, $hc\in\Omega_r$ has the phase and the edge column of $c$, so $T_m$ is in the same case at $c$ and at $hc$. Let $h$ have radius $R$, and let $T_m$ be forward at $c$, with $m\ge2+6R$. Put $j=2-m-3r\le-3(r+2R)$, $\hat c=c\sqcup_rK^-_j$, and $D=M_mc$, $D'=M_mhc$. By Lemma~\ref{c79:lem:hglue} and $(\mathrm H_1)$, $h{\hat c}=hc\sqcup_rK^-_j$. By Lemma~\ref{c79:lem:probe} for $c$ and $hc$, which have the same index $J$, $F^s{\hat c}=F^sD\sqcup_rK^+_{s+J}$ and $F^sh{\hat c}=F^sD'\sqcup_rK^+_{s+J}$ for large $s$. For $s$ so large that $q^+(s+J)\ge r+2R+1$, Lemma~\ref{c79:lem:hglue} gives $F^sh{\hat c}=hF^s{\hat c}=hF^sD\sqcup_{r+R}K^+_{s+J}$. By $(\mathrm H_1)$, $hF^sD\in\Omega_r$ has the phase of $F^sD$, which is that of $F^sD'$. So $F^sD'$ and $hF^sD=F^shD$ agree on $(-\infty,r]$ and lie in $\Omega_r$ with the same phase; hence they are equal, and $hD=D'$. For arbitrary $m$, Lemma~\ref{c79:lem:Tbasic}(c) gives $T_m=F^{6i}T_{m+6i}F^{-6i}$, and $h$ commutes with $F$. For backward moves use $\bar M_m=\rho M_{-m}\rho$ and $\rho h\rho\in\mathcal H$.
\end{proof}

\emph{Blocks.} For $c\in\Omega_r$ and $t\in\Z$ write $n(t)=t+\vartheta(c)\bmod 6$ and
\[
z_t=Y_c(t,r),\qquad y_t=Y_c(t,r-1),\qquad w_t=Y_c(t,r-2).
\]
A \emph{block} of $c$ is a set $\{t_1,t_1+1,\dots,t_1+5\}$ of times with $n(t_1)=1$; the time $t_j=t_1+j-1$ is at \emph{position} $j$, so that position $j$ is $n\equiv j\pmod 6$, and $(z_{t_1},z_{t_2})$ is the \emph{type} of the block. We also write $t_0=t_1-1$ and $t_7=t_6+1$, the last time of the previous block and the first of the next. For a configuration $D$ obtained from $c$ by moves, blocks, types and positions refer to $\vartheta(D)$, and $z^D,y^D,w^D$ are the columns $r,r-1,r-2$ of $Y_D$.

\begin{lemma}[blocks]\label{c79:lem:blocks}
Let $c\in\Omega_r$ and let $(p,q)$ be the type of a block of $c$.
\begin{enumerate}
\item[(a)] $(z_{t_1},\dots,z_{t_6})=(p,q,\neg p,\neg q,p,q)$.
\item[(b)] The next block has type $(p',q')$ with $p'=p\oplus(q\wedge\neg y_{t_6})$ and $q'=q\oplus(p'\wedge\neg y_{t_7})$. In particular $p'=1$ if $(p,q)=(1,0)$, and the $q$-value of the previous block is $q$ if $p=0$.
\item[(c)] $(p,q)\ne(0,0)$.
\item[(d)] The time $t$ is at position $j$ of a block of type $(p,q)$ of $c$ if and only if $-1-t$ is at position $7-j$ of a block of type $(q,p)$ of $\rho c$.
\end{enumerate}
\end{lemma}

\begin{proof}
(a) By Lemma~\ref{c79:lem:glue}(c), $z_{t+1}=\neg z_{t-1}$ at the positions $2,\dots,5$. (b) Lemma~\ref{c79:lem:glue}(c) at $t_6$ and $t_7$, where $n\equiv0,1$, gives $z_{t_7}=z_{t_5}\oplus(z_{t_6}\wedge\neg y_{t_6})$ and $z_{t_8}=z_{t_6}\oplus(z_{t_7}\wedge\neg y_{t_7})$, and $z_{t_5}=p$, $z_{t_6}=q$ by (a). For the last statement apply this to the previous block $(p_-,q_-)$: $q=q_-\oplus(p\wedge\neg y_{t_1})=q_-$ if $p=0$. (c) If $(p,q)=(0,0)$, then (b) gives $(p',q')=(0,0)$, and for the previous block $q_-=q\oplus(p\wedge\neg y_{t_1})=0$ and then $p_-=p\oplus(q_-\wedge\neg y_{t_0})=0$. So every block has type $(0,0)$, and by (a) $z_t=1$ exactly when $n(t)\in\{3,4\}$. By the last statement of Lemma~\ref{c79:lem:glue}(c), this is the $y$-bit at $r$ of $F^tV_{\omega}$, $\omega$ the phase of $c$; at $t=-1,0$ it says $c_r=(V_\omega)_r$, which contradicts $c\in\Omega_r$. (d) By Lemma~\ref{c79:lem:Tbasic}(c) and \eqref{eq:rho}, $\vartheta(\rho c)=2-\vartheta(c)$ and $Y_{\rho c}(t,i)=Y_c(-1-t,i)$; so the value of $n$ for $\rho c$ at $-1-t$ is $1-n(t)$, and the block is traversed in reverse order.
\end{proof}

\begin{lemma}[one move]\label{c79:lem:onemove}
Let $c\in\Omega_r$, let $n(m)\in\{0,2\}$ and $D=T_mc$.
\begin{enumerate}
\item[(a)] For $j\ge0$ the rows $m-1-j$ and $m+j$ of $Y_D$ and $Y_c$ agree on the columns $\le r-j$; in particular $z^D_t=z_t$ for $t\in\{m-1,m\}$, $y^D_t=y_t$ for $m-2\le t\le m+1$, and $w^D_t=w_t$ for $m-3\le t\le m+2$. For every $h\in\mathcal H$ and $j\ge0$, $Y_{hD}(t,r-j)\oplus Y_D(t,r-j)=Y_{hc}(t,r-j)\oplus Y_c(t,r-j)$ for $m-1-j\le t\le m+j$.
\item[(b)] Exactly one of the following holds.
\begin{itemize}
\item[$\mathrm F^-$:] $n(m)=2$ and $z_{m-1}=1$. Then $\vartheta(D)=\vartheta(c)-1$, and $m$ is at position $1$ of a block of $D$ of type $\bigl(z_m,\,1\oplus(z_m\wedge\neg y_m)\bigr)$.
\item[$\mathrm F^+$:] $n(m)=2$ and $z_{m-1}=0$. Then $\vartheta(D)=\vartheta(c)+2$, and $m$ is at position $4$ of a block of $D$ of type $(1,\neg z_m)$.
\item[$\mathrm B^-$:] $n(m)=0$ and $z_m=0$. Then $\vartheta(D)=\vartheta(c)-2$, and $m$ is at position $4$ of a block of $D$ of type $(\neg z_{m-1},1)$.
\item[$\mathrm B^+$:] $n(m)=0$ and $z_m=1$. Then $\vartheta(D)=\vartheta(c)+1$, $m$ is at position $1$ of a block of $D$ of type $(1,z_{m+1})$, and $m-1$ is at position $6$ of a block of $D$ of type $\bigl(1\oplus(z_{m-1}\wedge\neg y_{m-1}),\,z_{m-1}\bigr)$.
\end{itemize}
\end{enumerate}
\end{lemma}

\begin{proof}
(a) By Definition~\ref{c79:def:moves}, the rows $m-1$ and $m$ of $Y_D$ and $Y_c$ agree on the columns $\le r$; as $F^{\pm1}$ have radius $1$, the rows $m-1-j$ and $m+j$ agree on the columns $\le r-j$. The same holds for $hD$ and $hc$, since $hD=T_mhc$ by Lemma~\ref{c79:lem:commute}. (b) The four cases are those of Definition~\ref{c79:def:moves}, and $\vartheta(D)-\vartheta(c)$ is the change of phase. So the value of $n$ for $D$ at $m$ is $1,4,4,1$ in the four cases. In $\mathrm F^-$, $z^D_m=z_m$, and Lemma~\ref{c79:lem:glue}(c) for $D$ at $m$ gives $z^D_{m+1}=z^D_{m-1}\oplus(z^D_m\wedge\neg y^D_m)=1\oplus(z_m\wedge\neg y_m)$. In $\mathrm F^+$ and $\mathrm B^-$, $m-1$ and $m$ are at the positions $3$ and $4$, where by Lemma~\ref{c79:lem:blocks}(a) the values are the negations of the type; they are $z_{m-1}=0$ and $z_m$, respectively $z_{m-1}$ and $z_m=0$. In $\mathrm B^+$, $z^D_m=z_m=1$, and Lemma~\ref{c79:lem:glue}(c) at $m$, for $D$ (where $n\equiv1$) and for $c$ (where $n\equiv0$), gives $z^D_{m+1}=z_{m-1}\oplus\neg y_m=z_{m+1}$. Finally $m-1$ is at position $6$ for $D$, with $z^D_{m-1}=z_{m-1}$, and Lemma~\ref{c79:lem:glue}(c) for $D$ at $m-1$ gives $z^D_{m-2}=z^D_m\oplus(z^D_{m-1}\wedge\neg y^D_{m-1})=1\oplus(z_{m-1}\wedge\neg y_{m-1})$.
\end{proof}

\begin{proposition}[$(\mathrm H_2)$]\label{c79:prop:H2}
$(\mathrm H_2)$ holds.
\end{proposition}

\begin{proof}
For $h\in\mathcal H$ and $c\in\Omega_r$ put $g_t=g^c_t=Y_{hc}(t,r-1)\oplus Y_c(t,r-1)$. By $(\mathrm H_1)$ it suffices to show $g^c=0$, and $Y_{hc}$ and $Y_c$ agree on the columns $\ge r$ for every $c$. We use three rules, valid for every $h\in\mathcal H$ and $c\in\Omega_r$.
\begin{itemize}
\item[(G1)] If $z_t=1$ and $n(t)\in\{0,1\}$, then $g_t=0$: by Lemma~\ref{c79:lem:glue}(c) for $c$ and $hc$, which have the same phase and edge column, $z_{t+1}\oplus z_{t-1}=f(Y_c(t,r-1),1,1)=f(Y_{hc}(t,r-1),1,1)$, and $f(a,1,1)=\neg a$.
\item[(G2)] If $z_t=0$, then $g_{t-1}=g_{t+1}$: by \eqref{eq:rel} at $(t,r-1)$ and $f(a,b,0)=1$, $Y(t+1,r-1)=\neg Y(t-1,r-1)$ for both diagrams.
\item[(G3)] If $D=T_mc$, then $g^D_t=g^c_t$ for $m-2\le t\le m+1$, by Lemma~\ref{c79:lem:onemove}(a).
\end{itemize}
By Lemma~\ref{c79:lem:blocks}(c) every block of $c$ has type $(1,1)$, $(1,0)$ or $(0,1)$. We show that $g$ vanishes at every position of every block; $t_j$ denote the times of the block considered.

\emph{Step 1. Type $(1,1)$.} At $t_1$ and $t_6$ by (G1). The move at $t_2$ is $\mathrm F^-$, as $z_{t_1}=1$, and $t_2$ is at position $1$ of a block of $D=T_{t_2}c$ of type $(1,\cdot)$; so $g^D_{t_2}=0$ by (G1) for $D$, and $g_{t_2}=0$ by (G3). The move at $t_6$ is $\mathrm B^+$, as $z_{t_6}=1$, and $t_5$ is at position $6$ of a block of $D=T_{t_6}c$ whose $q$-value is $z_{t_5}=1$; (G1) for $D$ and (G3) give $g_{t_5}=0$. Finally (G2) at $t_3$ and $t_4$, where $z=0$, gives $g_{t_4}=g_{t_2}$ and $g_{t_3}=g_{t_5}$.

\emph{Step 2. Type $(0,1)$, positions $6$, $4$, $1$.} At $t_6$ by (G1), and $g_{t_4}=g_{t_6}$ by (G2) at $t_5$. The move at $t_2$ is $\mathrm F^+$, and $t_2$ is at position $4$, so $t_1$ at position $3$, of a block of $D=T_{t_2}c$ of type $(1,0)$. In $D$, (G1) at position $1$ and (G2) at position $2$ give $g^D=0$ at position $3$; so $g_{t_1}=0$ by (G3).

\emph{Step 3. Type $(1,0)$, positions $1$, $3$, $6$.} At $t_1$ by (G1), and $g_{t_3}=g_{t_1}$ by (G2) at $t_2$. The move at $t_6$ is $\mathrm B^-$, and $t_6$ is at position $4$ of a block of $D=T_{t_6}c$ of type $(\neg z_{t_5},1)=(0,1)$; by Step~2 for $D$ and (G3), $g_{t_6}=0$.

\emph{Step 4. Type $(1,0)$, positions $2$, $4$, $5$.} The move at $t_2$ is $\mathrm F^-$, and $t_2$ is at position $1$ of a block of $D=T_{t_2}c$ of type $(z_{t_2},1)=(0,1)$; by Step~2 for $D$ and (G3), $g_{t_2}=0$. Then $g_{t_4}=g_{t_2}$ by (G2) at $t_3$. By (G2) at $t_6$, $g_{t_5}=g_{t_7}$, and $t_7$ is at position $1$ of the next block, whose $p$-value is $1$ by Lemma~\ref{c79:lem:blocks}(b); so $g_{t_7}=0$ by (G1).

\emph{Step 5. Type $(0,1)$, positions $2$, $3$, $5$.} Steps 1--4 hold for every $h\in\mathcal H$ and every configuration. Apply Steps~3 and~4 to $\rho h\rho\in\mathcal H$ and $\rho c$: by Lemma~\ref{c79:lem:blocks}(d) the block becomes a block of type $(1,0)$ of $\rho c$, and by \eqref{eq:rho} the quantity $g$ for $\rho h\rho$ and $\rho c$ at $-1-t$ is $g^c_t$. Steps~3 and~4 are applied to $\rho c$ itself, with its own moves; translated back to $c$, the time $t'$ of $\rho c$ is the time $-1-t'$ of $c$, while a move of $\rho c$ at time $m'$ is the move of $c$ at time $-m'$ (Lemma~\ref{c79:lem:Tbasic}(c)), and the window $[m'-2,m'+1]$ of (G3) becomes $[-m'-2,-m'+1]$, the window of the move at $-m'$.
\end{proof}

\begin{proposition}[$(\mathrm H_3)$]\label{c79:prop:H3}
$(\mathrm H_3)$ holds.
\end{proposition}

\begin{proof}
For $h\in\mathcal H$ and $c\in\Omega_r$ put $g_t=g^c_t=Y_{hc}(t,r-2)\oplus Y_c(t,r-2)$. By $(\mathrm H_2)$ it suffices to show $g^c=0$, and $Y_{hc}$ and $Y_c$ agree on the columns $\ge r-1$ for every $c$. The rules are now, for every $h$ and $c$:
\begin{itemize}
\item[(K1)] if $y_t=z_t=1$, then $g_t=0$, by \eqref{eq:rel} at $(t,r-1)$ and $f(a,1,1)=\neg a$;
\item[(K2)] if $y_t=0$, then $g_{t-1}=g_{t+1}$, by \eqref{eq:rel} at $(t,r-2)$ and $f(a,b,0)=1$;
\item[(K3)] if $D=T_mc$, then $g^D_t=g^c_t$ for $m-3\le t\le m+2$, by Lemma~\ref{c79:lem:onemove}(a).
\end{itemize}
Moreover, by \eqref{eq:rel} at $(t,r)$, $(t,r-1)$, $(t,r-2)$ and $f(a,0,c)=\neg c$, $f(0,1,c)=1$, $f(a,b,0)=1$, if $y_t=0$ then
\[
\text{(Y)}\ \ y_{t+1}=y_{t-1}\oplus\neg z_t,\qquad \text{(Z)}\ \ z_{t+1}=\neg z_{t-1}\ \text{ if } z_t=1,\qquad \text{(W)}\ \ w_{t+1}=\neg w_{t-1}.
\]
Each of the following steps holds for every $h\in\mathcal H$, $c\in\Omega_r$ and $t$; by Lemma~\ref{c79:lem:blocks}(d) and \eqref{eq:rho}, applying a step to $\rho h\rho$ and $\rho c$ at $-1-t$ gives the mirror statement, in which position $j$ becomes $7-j$ with the type reversed. In it the times are reflected by $t\mapsto-1-t$ and the move times by $m\mapsto-m$ (Lemma~\ref{c79:lem:Tbasic}(c)): the time $t+s$ becomes $t-s$, and a move at time $t+s$ becomes a move at time $t+1-s$; the windows of (K3) correspond, since $t\in[m-3,m+2]$ if and only if $-1-t\in[-m-3,-m+2]$. This is what ``by $\rho$'' means below.

\emph{Step 1. If $y_t=1$, then $g_t=0$.} A move at a time $m$ with $t-1\le m\le t+2$ does not change $y_t$ and $g_t$ (Lemma~\ref{c79:lem:onemove}(a)). We show that at most four such moves lead to $D$ with $z^D_t=1$; then $g_t=g^D_t=0$ by (K3) and (K1) for $D$. Let $z_t=0$. By Lemma~\ref{c79:lem:blocks}(a), $t$ is at position $3$ or $4$ of a block of type $(1,1)$, at position $2$, $3$ or $6$ of type $(1,0)$, or at position $1$, $4$ or $5$ of type $(0,1)$. By $\rho$ it suffices to treat the positions $3$ of $(1,1)$ and $(1,0)$, $2$ of $(1,0)$ and $1$ of $(0,1)$.
\emph{Position $3$, type $(1,q)$.} The move at $t-1$ is $\mathrm F^-$, and $t$ is at position $2$ of a block of $D$ of type $\bigl(q,1\oplus(q\wedge\neg y_{t-1})\bigr)$. So $z^D_t=1$, unless $q=1$ and $y_{t-1}=0$, and then $t$ is at position $2$ of a block of $D$ of type $(1,0)$.
\emph{Position $2$, type $(1,0)$.} The move at $t$ is $\mathrm F^-$, and $t$ is at position $1$ of a block of $D$ of type $(0,1)$.
\emph{Position $1$, type $(0,1)$.} The move at $t+1$ is $\mathrm F^+$, and $t+1$ is at position $4$ of a block of $D$ of type $(1,0)$; so $t$ is at position $3$ of a block of type $(1,0)$, which is settled by one more move.

\emph{Step 2. If $y_t=0$ and $z_t=0$, then $g_t=0$.} By (Y), $y_{t+1}\ne y_{t-1}$; by $\rho$ we may assume $y_{t+1}=0$. Then $g_t=g_{t+2}$ by (K2) at $t+1$, and $y_{t+2}=\neg z_{t+1}$ by (Y) at $t+1$. If $z_{t+1}=0$, then $g_{t+2}=0$ by Step~1. If $z_{t+1}=1$, then (Z) and (Y) at $t+1$, $t+2$, $t+3$ give $z_{t+2}=1$, $y_{t+2}=y_{t+3}=0$, $z_{t+3}=0$ and $y_{t+4}=1$; so $g_{t+2}=g_{t+4}$ by (K2) at $t+3$, and $g_{t+4}=0$ by Step~1.

\emph{Step 3. If $y_t=0$, $z_t=1$ and $y_{t-1}y_{t+1}=0$, then $g_t=0$.} By (Y), $y_{t-1}=y_{t+1}=0$, and by (Z), $z_{t-1}\ne z_{t+1}$; by $\rho$ we may assume $z_{t+1}=0$. Then $g_t=g_{t+2}$ by (K2) at $t+1$, and $y_{t+2}=1$ by (Y) at $t+1$; so $g_{t+2}=0$ by Step~1.

\emph{Step 4. If $y_t=0$, $z_t=1$ and $y_{t-1}=y_{t+1}=1$, then $g_t=0$.} We find $D$, obtained from $c$ by one or two moves at times $m$ with $m-3\le t\le m+2$, with $z^D_t=0$ or $y^D_t=1$; then $g^D_t=0$ by Steps~1--3 for $D$, and $g_t=g^D_t$ by (K3). By (Z) and (W), $z_{t-1}\ne z_{t+1}$ and $w_{t-1}\ne w_{t+1}$; by $\rho$ we may assume $z_{t-1}=0$ and $z_{t+1}=1$. By Lemma~\ref{c79:lem:blocks}(a), $t$ is then at position $1$ of a block of type $(1,1)$ (and $z_{t-1}=0$ is the $q$-value of the previous block), at position $5$ of type $(1,1)$, at position $2$ of type $(0,1)$, at position $6$ of type $(0,1)$ (and $z_{t+1}=1$ is the $p$-value of the next block), or at position $4$ of type $(1,0)$.
\emph{Position $1$, type $(1,1)$.} The previous block has type $(1,0)$, by Lemma~\ref{c79:lem:blocks}(c), so $z_{t-2}=1$. The move at $t-1$ is $\mathrm B^-$, and $t-1$ is at position $4$ of a block of $D$ of type $(\neg z_{t-2},1)=(0,1)$; $t$ is at position $5$, so $z^D_t=0$.
\emph{Position $5$, type $(1,1)$.} By Lemma~\ref{c79:lem:blocks}(b) the next block has $p$-value $1\oplus\neg y_{t+1}=1$, so the move at $t+3$ is $\mathrm F^-$. Then $t+3$ is at position $1$ of a block of $D$, so $t+2$ is at position $6$ of a block of $D$ whose $q$-value is $z^D_{t+2}=z_{t+2}=1$, and $t$, at position $4$ of that block, has $z^D_t=0$.
\emph{Position $2$, type $(0,1)$.} By Lemma~\ref{c79:lem:blocks}(b) the previous block has $q$-value $1$, so the move at $t-2$ is $\mathrm B^+$, and $t-2$ is at position $1$ of a block of $D$ of type $(1,z_{t-1})=(1,0)$; $t$ is at position $3$, so $z^D_t=0$.
\emph{Position $6$, type $(0,1)$.} The move at $t$ is $\mathrm B^+$, and $t-1$ is at position $6$ of a block of $D$ of type $\bigl(1\oplus(z_{t-1}\wedge\neg y_{t-1}),z_{t-1}\bigr)=(1,0)$. In $D$ the move at $t-1$ is $\mathrm B^-$, and $t-1$ is at position $4$ of a block of $D'=T_{t-1}D$ of type $(\neg z^D_{t-2},1)=(0,1)$; $t$ is at position $5$, so $z^{D'}_t=0$.
\emph{Position $4$, type $(1,0)$.} The move at $t-2$ is $\mathrm F^-$, and $t-2$ is at position $1$ of a block of $D$ of type $(z_{t-2},1)=(0,1)$; so $z^D_{t-1}=z^D_t=1$ and $z^D_{t+1}=0$. By \eqref{eq:rel} at $(t-1,r-1)$ for $D$ and for $c$, with $y^D_{t-2}=y_{t-2}$, $y^D_{t-1}=y_{t-1}=1$ and $w^D_{t-1}=w_{t-1}$,
\[
y^D_t=y_{t-2}\oplus f(w_{t-1},1,1)=y_{t-2}\oplus\neg w_{t-1},\qquad 0=y_t=y_{t-2}\oplus f(w_{t-1},1,0)=y_{t-2}\oplus1,
\]
so $y^D_t=w_{t-1}$. If $w_{t-1}=1$, then $y^D_t=1$. If $w_{t-1}=0$, then $y^D_t=0$, $z^D_t=1$, $y^D_{t-1}=1$, and $w^D_{t+1}=\neg w^D_{t-1}=1$ by (W) for $D$: so $(D,t)$ is in the situation of Step~4 with $z^D_{t-1}=1$, $z^D_{t+1}=0$, at position $3$ of a block of type $(0,1)$, and with $w^D_{t+1}=1$. This is the mirror of the case just treated with $w_{t-1}=1$; so one more move, at $t+3$, gives $D'$ with $y^{D'}_t=1$.

By Steps~1--4, $g_t=0$ for every $t$.
\end{proof}


\subsection{Controllability}\label{c79:sec:Fk}

For $k\ge1$ and $c\in\Omega_r$ let $\mathcal O_k(c)$ be the set of configurations $T_{m_n}\cdots T_{m_1}c$ with $n\ge0$ and $|m_1|,\dots,|m_n|\le k$. As in Section~\ref{cl:sec:ctrl}, $c_{r-k+1}$ is the tip and $c_{r-k}$ the exterior cell of $c$ at level $k$.

The orbits at level $3$ are almost those of Theorem~\ref{c79:thm:Fk}. We determine them in Proposition~\ref{c79:prop:F3}, from an explicit description of the moves at the times $|m|\le3$.

\emph{Notation at level $3$.} For $c\in\Omega_r$ write $a_t=Y_c(t,r-3)$, in addition to $z_t$, $y_t$, $w_t$; so
\[
c_{r-3}=(a_{-1},a_0),\qquad c_{r-2}=(w_{-1},w_0),\qquad c_{r-1}=(y_{-1},y_0),\qquad c_r=(z_{-1},z_0).
\]
A cell $(x,y)$ is also written as the string $xy$, and $c$ is abbreviated to
\[
\cfg\vartheta\tau\mu\varepsilon\qquad\text{with}\qquad\vartheta=\vartheta(c),\quad\tau=c_{r-2},\quad\mu=c_{r-1},\quad\varepsilon=c_r.
\]
This determines $c$ when its cells on $(-\infty,r-3]$ are given, and by Lemma~\ref{c79:lem:Tbasic}(b) no move at a time $|m|\le3$ changes these cells. By the last statement of Lemma~\ref{c79:lem:glue}(c), the cell at $r$ of the vacuum of $c$ is $00$, $00$, $00$, $01$, $11$, $10$ for $\vartheta=0,\dots,5$, and $\varepsilon$ can be any other cell; so there are eighteen pairs $(\vartheta,\varepsilon)$. For bits $u,v$ the bracket $[uv\ne01]$ is $0$ if $(u,v)=(0,1)$ and $1$ otherwise; by \eqref{c79:eq:f},
\begin{equation}\label{c79:eq:br}
f(u,v,0)\oplus f(u,v,1)=[uv\ne01],\qquad f(u,v,1)=\neg u\wedge v .
\end{equation}
For $c,D\in\Omega_r$ we write $c\mv mD$ if $D=T_mc\ne c$, and we say that $c$ \emph{reaches} $D$ if $D\in\mathcal O_3(c)$; this relation is transitive.

\emph{Mirror images.} By Lemma~\ref{c79:lem:Tbasic}(c), $\rho$ maps $\Omega_r$ onto itself, it exchanges the two bits of every cell, $\vartheta(\rho c)=2-\vartheta(c)$, and $\bar M_m=\rho M_{-m}\rho$, hence also $M_m=\rho\bar M_{-m}\rho$. Moreover $-m+\vartheta(\rho c)=2-(m+\vartheta(c))$, so $T_{-m}$ is backward, forward or the identity at $\rho c$ exactly when $T_m$ is forward, backward or the identity at $c$. So $T_{-m}=\rho T_m\rho$ for every $m$, and $c\mv mD$ if and only if $\rho c\mv{-m}\rho D$. Therefore every statement below that holds for all $c\in\Omega_r$ has a \emph{mirror image}, which also holds for all $c\in\Omega_r$: in it $\vartheta$ is replaced by $2-\vartheta$, the two bits of every cell are exchanged, including those of $c_{r-3}$, and $T_m$ is replaced by $T_{-m}$. We write $\bar\tau$ for the cell $\tau$ with its two bits exchanged.

\begin{lemma}[the moves at the times $0,\dots,3$]\label{c79:lem:M3}
Let $c\in\Omega_r$, let $m\in\{0,1,2,3\}$ with $n(m)\in\{0,2\}$, and let $D=T_mc$. Write $\vartheta'$, $z'_t$, $y'_t$, $w'_t$ for the quantities $\vartheta$, $z_t$, $y_t$, $w_t$ of $D$, and put
\[
y_1=y_{-1}\oplus f(w_0,y_0,z_0),\qquad w_1=w_{-1}\oplus f(a_0,w_0,y_0),
\]
which are $Y_c(1,r-1)$ and $Y_c(1,r-2)$ by \eqref{eq:rel}.
\begin{enumerate}
\item[(a)] The value $\vartheta'$, the edge cell $D_r=(z'_{-1},z'_0)$ and, for $m=3$, the value $z'_1$ are given by the following table. In it $z_1=z_{-1}\oplus(z_0\wedge\neg y_0)$ if $\vartheta=0$, and $z_1=\neg z_{-1}$ if $\vartheta\in\{3,5\}$; $z_2=z_0\oplus(z_1\wedge\neg y_1)$ if $\vartheta=5$; and $\pi=1\oplus(\neg z_0\wedge y_0)$.
\[
\begin{array}{cclcll}
m&\vartheta&\text{case}&\vartheta'&D_r&\\\hline
0&2&z_{-1}=0&4&c_r&\\
0&2&z_{-1}=1&1&c_r&\\
0&0&z_0=0&4&c_r&\\
0&0&z_0=1&1&c_r&\\
1&1&z_0=0&3&00&\\
1&1&z_0=1&0&c_r&\\
1&5&z_{-1}=1&3&c_r&\\
1&5&z_{-1}=0&0&\bigl(1\oplus(z_0\wedge\neg y_0),\,z_0\bigr)&\\
2&0&z_1=0&2&10&\\
2&0&z_1=1&5&(0,z_0)&\\
2&4&z_0=1&2&c_r&\\
2&4&z_0=0&5&\bigl(z_{-1},\,1\oplus(\neg z_{-1}\wedge y_{-1})\bigr)&\\
3&5&z_2=0&1&(\neg y'_0,1)&z'_1=0,\ z_1=1\\
3&5&z_2=1&4&(z_{-1},0)&z'_1=z_1\\
3&3&z_{-1}=0&1&10&z'_1=z_1\\
3&3&z_{-1}=1&4&(\neg\pi,z_0)&z'_1=\pi,\ z_1=0
\end{array}
\]
\item[(b)] $D$ agrees with $c$ on $(-\infty,r-m]$. If $m=2$, or if $m=3$ and $z'_1=z_1$, then $D_{r-2}=c_{r-2}$, $y'_0=y_0$ and
\[
y'_{-1}=y_{-1}\oplus\bigl([w_0y_0\ne01]\wedge(z_0\oplus z'_0)\bigr).
\]
If $m=3$ and $z'_1\ne z_1$, then, with $\delta=[w_1y_1\ne01]$,
\[
y'_0=y_0\oplus\delta,\qquad w'_0=w_0,\qquad w'_{-1}=w_{-1}\oplus\bigl(\delta\wedge[a_0w_0\ne01]\bigr),\qquad y'_{-1}=y_1\oplus f(w_0,y'_0,z'_0).
\]
\end{enumerate}
\end{lemma}

\begin{proof}
(b) The first statement is Lemma~\ref{c79:lem:Tbasic}(b). By Lemma~\ref{c79:lem:onemove}(a), for $j\ge0$ the row $m-1-j$ of $Y_D$ agrees with that of $Y_c$ on the columns $\le r-j$, and the row $m$ on the columns $\le r$. Let $m=2$. Then $y'_0=y_0$, $y'_1=y_1$ and $w'_0=w_0$, and \eqref{eq:rel} at $(0,r-1)$, for $D$ and for $c$, gives $y'_{-1}\oplus y_{-1}=f(w_0,y_0,z'_0)\oplus f(w_0,y_0,z_0)$, which is $[w_0y_0\ne01]\wedge(z_0\oplus z'_0)$ by \eqref{c79:eq:br}. Let $m=3$. Then $y'_1=y_1$, $y'_2=y_2$, $w'_0=w_0$, $w'_1=w_1$ and $Y_D(0,r-3)=a_0$. By \eqref{eq:rel} at $(1,r-1)$, $(0,r-2)$ and $(0,r-1)$, for $D$ and for $c$,
\[
y'_0\oplus y_0=[w_1y_1\ne01]\wedge(z_1\oplus z'_1),\quad w'_{-1}\oplus w_{-1}=[a_0w_0\ne01]\wedge(y_0\oplus y'_0),\quad y'_{-1}=y_1\oplus f(w_0,y'_0,z'_0).
\]
If $z'_1=z_1$, then $y'_0=y_0$ and $w'_{-1}=w_{-1}$, and $y'_{-1}\oplus y_{-1}$ is found as for $m=2$. If $z'_1\ne z_1$, these are the formulas of (b).

(a) For $m=0$ this is Definition~\ref{c79:def:moves}. Let $m\ge1$. The value of $\vartheta'$ is given by Lemma~\ref{c79:lem:onemove}(b), and the case is decided by $z_{m-1}$ if $n(m)=2$ and by $z_m$ if $n(m)=0$. By Lemma~\ref{c79:lem:blocks}(a), at the positions $1,\dots,6$ of a block of type $(p,q)$ the edge column is $p,q,\neg p,\neg q,p,q$, for $c$ and for $D$. For $c$ this gives: if $\vartheta=5$, the times $-1,0,1$ are at the positions $4,5,6$, so $z_1=\neg z_{-1}$; if $\vartheta=4$, the times $-1,\dots,2$ are at the positions $3,\dots,6$, so $z_1=\neg z_{-1}$ and $z_2=\neg z_0$; if $\vartheta=3$, the times $-1,\dots,3$ are at the positions $2,\dots,6$, so $z_1=\neg z_{-1}$, $z_2=\neg z_0$ and $z_3=z_{-1}$. By Lemma~\ref{c79:lem:glue}(c), $z_1=z_{-1}\oplus(z_0\wedge\neg y_0)$ if $\vartheta=0$, and $z_2=z_0\oplus(z_1\wedge\neg y_1)$ if $\vartheta=5$. So the cases of the table are those of Lemma~\ref{c79:lem:onemove}(b), in the order $\mathrm F^+$, $\mathrm F^-$ for $(m,\vartheta)=(1,1),(2,0),(3,5)$ and $\mathrm B^-$, $\mathrm B^+$ for $(m,\vartheta)=(1,5),(2,4),(3,3)$. We determine $z'_t$ for $m-4\le t\le m-1$ in the four cases; recall $z'_{m-1}=z_{m-1}$ and $z'_m=z_m$.

$\mathrm F^+$. Here $z_{m-1}=0$, so $z_m=1$ by Lemma~\ref{c79:lem:blocks}(c), and $m$ is at position $4$ of a block of $D$ of type $(1,0)$; so $(z'_{m-3},\dots,z'_m)=(1,0,0,1)$. This gives $D_r=00$ for $m=1$ and $D_r=10$ for $m=2$. For $m=3$ it gives $z'_0=1$ and $z'_1=0$, and Lemma~\ref{c79:lem:glue}(c) for $D$ at $t=0$, where the value of $n$ for $D$ is $\vartheta'=1$, gives $z'_1=z'_{-1}\oplus(z'_0\wedge\neg y'_0)$, so $z'_{-1}=\neg y'_0$. Moreover $z_1=1$: the block of $c$ that begins at time $2$ has type $(0,1)$, so by Lemma~\ref{c79:lem:blocks}(b) the previous block, which ends at time $1$, has $q$-value $1$.

$\mathrm F^-$. Here $z_{m-1}=1$, and $m$ is at position $1$ of a block of $D$, so $m-1$ is at position $6$ of the previous block of $D$, whose $q$-value is $z'_{m-1}=1$. Lemma~\ref{c79:lem:glue}(c) at $t=m-1$, for $D$ (where $n=0$) and for $c$ (where $n=1$), gives $z'_m=z'_{m-2}\oplus(z'_{m-1}\wedge\neg y'_{m-1})$ and the same relation without primes; as $y'_{m-1}=y_{m-1}$, it follows that $z'_{m-2}=z_{m-2}$. So that block of $D$ has type $(z_{m-2},1)$, and $(z'_{m-4},z'_{m-3},z'_{m-2})=(\neg z_{m-2},0,z_{m-2})$. This gives $D_r=c_r$ for $m=1$, $D_r=(0,z_0)$ for $m=2$, and $D_r=(\neg z_1,0)=(z_{-1},0)$, $z'_1=z_1$ for $m=3$.

$\mathrm B^-$. Here $z_m=0$, so $z_{m-1}=1$ by Lemma~\ref{c79:lem:blocks}(c), and $m$ is at position $4$ of a block of $D$ of type $(0,1)$; so $(z'_{m-3},\dots,z'_m)=(0,1,1,0)$, and $z'_{m-4}=1$, the $q$-value of the previous block of $D$ (Lemma~\ref{c79:lem:blocks}(b)). For $c$, $m$ is at position $6$ of a block of type $(1,0)$, so $(z_{m-4},\dots,z_m)=(0,0,1,1,0)$. Hence $D_r=c_r$ for $m=1,2$, and $D_r=10$, $z'_1=z_1$ for $m=3$.

$\mathrm B^+$. Here $z_m=1$, and $m-1$ is at position $6$ of a block of $D$ of type $(\pi_m,z_{m-1})$ with $\pi_m=1\oplus(z_{m-1}\wedge\neg y_{m-1})$; so $(z'_{m-4},\dots,z'_{m-1})=(\neg\pi_m,\neg z_{m-1},\pi_m,z_{m-1})$. For $m=1$, $D_r=(\pi_1,z_0)$. For $m=2$, $D_r=(\neg z_1,\pi_2)=(z_{-1},\pi_2)$; here $z_0=\neg z_2=0$, so $y_1=y_{-1}\oplus f(w_0,y_0,0)=\neg y_{-1}$ and $\pi_2=1\oplus(\neg z_{-1}\wedge y_{-1})$. For $m=3$, $D_r=(\neg\pi_3,\neg z_2)=(\neg\pi_3,z_0)$ and $z'_1=\pi_3$; here $z_1=\neg z_{-1}=\neg z_3=0$, so $y_2=y_0\oplus f(w_1,y_1,0)=\neg y_0$ and $\pi_3=1\oplus(\neg z_0\wedge y_0)=\pi$.
\end{proof}

Let $\mathcal E$ be the set of configurations $b\in\Omega_r$, for any $r$, with
\[
b_{r-2}=(0,0),\qquad b_{r-1}=(1,1),\qquad (\vartheta(b),b_r)\in\bigl\{(0,(1,1)),\,(1,(1,1)),\,(2,(1,1)),\,(3,(1,0)),\,(5,(0,1))\bigr\}.
\]

\begin{proposition}[$(\mathrm F^\ast_3)$]\label{c79:prop:F3}
Let $c\in\Omega_r$. Then $\mathcal O_3(c)$ contains every $c'\in\Omega_r$, of any phase, with $c'\notin\mathcal E$, $c'_i=c_i$ for $i\le r-3$, and, if $c_{r-3}=(0,0)$, $c'_{r-2}=(1,1)$ if and only if $c_{r-2}=(1,1)$.
\end{proposition}

\begin{proof}
The proof has four steps. Twelve of the eighteen pairs $(\vartheta,\varepsilon)$ occur in sixteen \emph{sectors} of twelve configurations each, and any two configurations of a sector reach each other (Step~1). A configuration with one of the other six pairs reaches a sector (Step~2) and, unless it lies in $\mathcal E$, is reached from one (Step~3). Moves from sector to sector, mainly the moves at the times $\pm3$ applied to the configurations with $\vartheta\in\{3,5\}$ and edge cell $11$, join the sectors as far as the condition on $(1,1)$ allows (Step~4).

All configurations in this proof lie in $\Omega_r$ and agree with $c$ on $(-\infty,r-3]$. Every arrow stated below is proved for arbitrary cells on $(-\infty,r-3]$; so its mirror image holds as well, for arbitrary cells on $(-\infty,r-3]$. In a configuration $\cfg\vartheta\tau\nu\varepsilon$ we always write $\tau=(w_{-1},w_0)$, $\nu=(y_{-1},y_0)$, $\varepsilon=(z_{-1},z_0)$, and $(a_{-1},a_0)=c_{r-3}$; the middle cell is called $\nu$ or $\mu$. For cells $\tau=(w_{-1},w_0)$ and $\nu=(y_{-1},y_0)$ put
\[
\nu^x=\bigl(y_{-1}\oplus[w_0y_0\ne01],\ y_0\bigr),\qquad \nu^y=\bigl(y_{-1},\ y_0\oplus[w_{-1}y_{-1}\ne01]\bigr).
\]
These cells depend on $\tau$, which is always the cell written to the left of $\nu$. Since $\nu^x$ has the $y$-bit of $\nu$ and $\nu^y$ has its $x$-bit, $(\nu^x)^x=\nu$ and $(\nu^y)^y=\nu$. The mirror image of the operation $\nu\mapsto\nu^x$ is $\nu\mapsto\nu^y$.

\emph{The arrows of Steps~1--3.} Table~\ref{c79:tab:arrows} lists the arrows that are used in Steps~1--3. Each of them holds for all cells $\tau$, $\nu$ that satisfy the condition in its third column, and is read off Lemma~\ref{c79:lem:M3}: the fourth column names the row of the table of the lemma, by $(m,\vartheta)$ and the case, and the last column evaluates this row. An arrow at a time $m<0$ is the mirror image of the arrow at the time $-m$ that is displayed in the last column, and the fourth column names the row for the displayed arrow. The cells $c_{r-2}$ and $c_{r-1}$ are found from Lemma~\ref{c79:lem:M3}(b). If $m\le1$, they do not change. If $m=2$, or if $m=3$ and $z'_1=z_1$, then $\tau$ does not change, and $\nu$ is replaced by $\nu^x$ if $z_0$ changes and does not change otherwise. In the two arrows with $m=3$ and $z'_1\ne z_1$, which are \eqref{c79:eq:j11} and \eqref{c79:eq:j12}, the last column evaluates the second set of formulas of Lemma~\ref{c79:lem:M3}(b).

\begin{table}[tp]
\centering\footnotesize
\setlength{\tabcolsep}{3.5pt}
\begin{tabular}{@{}llll>{\raggedright\arraybackslash}p{7.3cm}@{}}
&arrow&condition&row&evaluation\\\hline
\arw{c79:eq:s1}&$\cfg4\tau\nu{10}\mv{-2}\cfg0\tau\nu{10}$&&$(2,4)$, $z_0=1$&mirror image of $\cfg4{\bar\tau}{\bar\nu}{01}\mv2\cfg2{\bar\tau}{\bar\nu}{01}$, where $D_r=c_r$\\
&$\cfg0\tau\nu{10}\mv{2}\cfg5\tau\nu{00}$&&$(2,0)$, $z_1=1$&$z_1=z_{-1}\oplus(z_0\wedge\neg y_0)=1$; $D_r=(0,z_0)=00$, and $z_0$ does not change\\
&$\cfg5\tau\nu{00}\mv{-3}\cfg1\tau\nu{01}$&&$(3,3)$, $z_{-1}=0$&mirror image of $\cfg3{\bar\tau}{\bar\nu}{00}\mv3\cfg1{\bar\tau}{\bar\nu}{10}$, where $D_r=10$, $z'_1=z_1$, and $z_0$ does not change\\
&$\cfg1\tau\nu{01}\mv{1}\cfg0\tau\nu{01}$&&$(1,1)$, $z_0=1$&$D_r=c_r$\\
&$\cfg0\tau\nu{01}\mv{0}\cfg1\tau\nu{01}$&&$(0,0)$, $z_0=1$&$D_r=c_r$\\
&$\cfg1\tau\nu{01}\mv{-1}\cfg5\tau\nu{00}$&&$(1,1)$, $z_0=0$&mirror image of $\cfg1{\bar\tau}{\bar\nu}{10}\mv1\cfg3{\bar\tau}{\bar\nu}{00}$, where $D_r=00$\\
&$\cfg5\tau\nu{00}\mv{1}\cfg0\tau\nu{10}$&&$(1,5)$, $z_{-1}=0$&$D_r=(1\oplus(0\wedge\neg y_0),0)=10$\\
&$\cfg0\tau\nu{10}\mv{0}\cfg4\tau\nu{10}$&&$(0,0)$, $z_0=0$&$D_r=c_r$\\[1pt]
\arw{c79:eq:s2}&$\cfg4\tau\nu{10}\mv{2}\cfg5\tau{\nu^x}{11}$&&$(2,4)$, $z_0=0$&$D_r=(1,1\oplus(\neg1\wedge y_{-1}))=11$, and $z_0$ changes\\
&$\cfg5\tau\nu{11}\mv{3}\cfg4\tau{\nu^x}{10}$&&$(3,5)$, $z_2=1$&$z_1=\neg z_{-1}=0$ and $z_2=z_0\oplus(z_1\wedge\neg y_1)=1$; $D_r=(z_{-1},0)=10$, $z'_1=z_1$, and $z_0$ changes\\
&$\cfg5\tau\nu{11}\mv{1}\cfg3\tau\nu{11}$&&$(1,5)$, $z_{-1}=1$&$D_r=c_r$\\
&$\cfg3\tau\nu{11}\mv{-1}\cfg5\tau\nu{11}$&&$(1,5)$, $z_{-1}=1$&mirror image of $\cfg5{\bar\tau}{\bar\nu}{11}\mv1\cfg3{\bar\tau}{\bar\nu}{11}$\\[1pt]
\arw{c79:eq:i3}&$\cfg0\tau\nu{11}\mv{0}\cfg1\tau\nu{11}$&&$(0,0)$, $z_0=1$&$D_r=c_r$\\
&$\cfg1\tau\nu{11}\mv{-1}\cfg2\tau\nu{11}$&&$(1,1)$, $z_0=1$&mirror image of $\cfg1{\bar\tau}{\bar\nu}{11}\mv1\cfg0{\bar\tau}{\bar\nu}{11}$\\
&$\cfg2\tau\nu{11}\mv{0}\cfg1\tau\nu{11}$&&$(0,2)$, $z_{-1}=1$&$D_r=c_r$\\
&$\cfg1\tau\nu{11}\mv{1}\cfg0\tau\nu{11}$&&$(1,1)$, $z_0=1$&$D_r=c_r$\\[1pt]
\arw{c79:eq:i4}&$\cfg0\tau\nu{11}\mv{2}\cfg2\tau{\nu^x}{10}$&$y_0=0$&$(2,0)$, $z_1=0$&$z_1=1\oplus(1\wedge\neg y_0)=y_0$; $D_r=10$, and $z_0$ changes\\
&$\cfg0\tau\nu{11}\mv{2}\cfg5\tau\nu{01}$&$y_0=1$&$(2,0)$, $z_1=1$&$z_1=y_0$; $D_r=(0,z_0)=01$\\[1pt]
\arw{c79:eq:i5}&$\cfg5\tau\nu{01}\mv{1}\cfg0\tau\nu{\varepsilon'}$&&$(1,5)$, $z_{-1}=0$&$\varepsilon'=D_r=(1\oplus(1\wedge\neg y_0),1)=(y_0,1)$\\[1pt]
\arw{c79:eq:i6}&$\cfg3\tau\nu{10}\mv{3}\cfg4\tau\nu{10}$&$y_0=1$&$(3,3)$, $z_{-1}=1$&$\pi=1\oplus(\neg z_0\wedge y_0)=0$; $D_r=(\neg\pi,z_0)=10$, $z'_1=\pi=0=z_1$, and $z_0$ does not change\\[1pt]
\arw{c79:eq:i7}&$\cfg4\tau\nu{00}\mv{2}\cfg5\tau\nu{00}$&$y_{-1}=1$&$(2,4)$, $z_0=0$&$D_r=(0,1\oplus(\neg0\wedge y_{-1}))=00$, and $z_0$ does not change\\
&$\cfg4\tau\nu{00}\mv{2}\cfg5\tau{\nu^x}{01}$&$y_{-1}=0$&$(2,4)$, $z_0=0$&$D_r=(0,1\oplus(\neg0\wedge y_{-1}))=01$, and $z_0$ changes\\[1pt]
\arw{c79:eq:j8}&$\cfg5\tau\nu{00}\mv{3}\cfg4\tau\nu{00}$&$y_{-1}=1$&$(3,5)$, $z_2=1$&$z_1=\neg z_{-1}=1$, $y_1=y_{-1}\oplus f(w_0,y_0,0)=\neg y_{-1}=0$ and $z_2=z_0\oplus\neg y_1=1$; $D_r=(z_{-1},0)=00$, $z'_1=z_1$, and $z_0$ does not change\\[1pt]
\arw{c79:eq:j9}&$\cfg5\tau{10}{01}\mv{3}\cfg4\tau{00}{00}$&&$(3,5)$, $z_2=1$&$z_1=1$, $y_1=1\oplus f(w_0,0,1)=1$ and $z_2=z_0\oplus\neg y_1=1$; $D_r=00$, $z'_1=z_1$, $z_0$ changes, and $\nu^x=(1\oplus[w_0\,0\ne01],0)=00$\\[1pt]
\arw{c79:eq:j10}&$\cfg0\tau\nu{01}\mv{2}\cfg5\tau\nu{01}$&$y_0=0$&$(2,0)$, $z_1=1$&$z_1=0\oplus(1\wedge\neg y_0)=1$; $D_r=(0,z_0)=01$\\[1pt]
\arw{c79:eq:j11}&$\cfg5\tau{00}{00}\mv{3}\cfg1\tau{10}{11}$&$w_{-1}=1$&$(3,5)$, $z_2=0$&$z_1=1$, $y_1=f(w_0,0,0)=1$ and $z_2=z_0\oplus\neg y_1=0$; so $z'_1=0\ne z_1$. $w_1=w_{-1}\oplus f(a_0,w_0,0)=0$ and $\delta=[01\ne01]=0$; hence $y'_0=0$, $w'_{-1}=w_{-1}$, $D_r=(\neg y'_0,1)=11$ and $y'_{-1}=y_1\oplus f(w_0,0,1)=1$\\[1pt]
\arw{c79:eq:j12}&$\cfg5\tau\nu{01}\mv{3}\cfg1{\tau'}{00}{11}$&$\nu=(\neg w_0,1)$&$(3,5)$, $z_2=0$&$z_1=1$, $y_1=\neg w_0\oplus f(w_0,1,1)=0$ and $z_2=z_0\oplus\neg y_1=0$; so $z'_1=0\ne z_1$. $\delta=[w_1\,0\ne01]=1$; hence $y'_0=0$, $\tau'=(w_{-1}\oplus[a_0w_0\ne01],\,w_0)$, $D_r=(\neg y'_0,1)=11$ and $y'_{-1}=y_1\oplus f(w_0,0,1)=0$\\\hline
\end{tabular}
\smallskip
\caption{The arrows of Steps~1--3 in the proof of Proposition~\ref{c79:prop:F3}. The eight arrows of \eqref{c79:eq:s1} form a closed chain, and so do the four arrows of \eqref{c79:eq:i3}.}\label{c79:tab:arrows}
\end{table}

\emph{Step 1: sectors.} For cells $\tau,\nu$ let
\begin{align*}
\mathsf H(\tau,\nu)&=\{\cfg3\tau\nu{11},\ \cfg5\tau\nu{11}\},\\
\mathsf A(\tau,\nu)&=\{\cfg4\tau\nu{10},\ \cfg0\tau\nu{10},\ \cfg5\tau\nu{00},\ \cfg1\tau\nu{01},\ \cfg0\tau\nu{01}\},\\
\mathsf D(\tau,\nu)&=\{\cfg4\tau\nu{01},\ \cfg2\tau\nu{01},\ \cfg3\tau\nu{00},\ \cfg1\tau\nu{10},\ \cfg2\tau\nu{10}\}.
\end{align*}
The mirror image of $\mathsf A(\tau,\nu)$ is $\mathsf D(\bar\tau,\bar\nu)$, and that of $\mathsf H(\tau,\nu)$ is $\mathsf H(\bar\tau,\bar\nu)$.

The arrows of \eqref{c79:eq:s1} form a closed chain through the five configurations of $\mathsf A(\tau,\nu)$; so these reach each other, and by the mirror image so do those of $\mathsf D(\tau,\nu)$. By \eqref{c79:eq:s2}, applied to $\nu=\mu^x$ and to $\nu=\mu$, the configurations of $\mathsf H(\tau,\mu)$ and of $\mathsf A(\tau,\mu^x)$ reach each other, and by the mirror image so do those of $\mathsf H(\tau,\mu)$ and of $\mathsf D(\tau,\mu^y)$. Hence any two of the twelve configurations of the \emph{sector}
\[
\Sigma(\tau,\mu)=\mathsf H(\tau,\mu)\cup\mathsf A(\tau,\mu^x)\cup\mathsf D(\tau,\mu^y)
\]
reach each other. The mirror image of $\Sigma(\tau,\mu)$ is $\Sigma(\bar\tau,\bar\mu)$. Twelve of the eighteen pairs $(\vartheta,\varepsilon)$ occur in $\mathsf H$, $\mathsf A$, $\mathsf D$, and a configuration $\cfg\vartheta\tau\nu\varepsilon$ with one of them lies in a sector: $\mathsf H(\tau,\nu)\subseteq\Sigma(\tau,\nu)$, $\mathsf A(\tau,\nu)\subseteq\Sigma(\tau,\nu^x)$ and $\mathsf D(\tau,\nu)\subseteq\Sigma(\tau,\nu^y)$. The other six pairs are $(0,11)$, $(1,11)$, $(2,11)$, $(5,01)$, $(3,10)$ and $(4,00)$; we call a configuration with one of them \emph{irregular}, and put $\mathsf C(\tau,\nu)=\{\cfg0\tau\nu{11},\cfg1\tau\nu{11},\cfg2\tau\nu{11}\}$. Thus $\mathcal E$ consists of $\mathsf C(00,11)$, $\cfg5{00}{11}{01}$ and $\cfg3{00}{11}{10}$.

\emph{Step 2: every irregular configuration reaches a configuration in a sector.} Let an irregular configuration be given.
\begin{itemize}
\item $\cfg4\tau\nu{00}$. If $y_{-1}=1$, \eqref{c79:eq:i7} leads to $\cfg5\tau\nu{00}\in\mathsf A(\tau,\nu)$. If $y_0=1$, the mirror image of \eqref{c79:eq:i7} leads to $\cfg3\tau\nu{00}\in\mathsf D(\tau,\nu)$. If $\nu=00$, then $\nu^x=10$, and \eqref{c79:eq:i7} and \eqref{c79:eq:i5} lead to $\cfg5\tau{10}{01}$ and to $\cfg0\tau{10}{01}\in\mathsf A(\tau,10)$.
\item $\cfg5\tau\nu{01}$. If $y_0=0$, \eqref{c79:eq:i5} leads to $\cfg0\tau\nu{01}\in\mathsf A(\tau,\nu)$. If $y_{-1}=1$, the mirror image of \eqref{c79:eq:i6} leads to $\cfg4\tau\nu{01}\in\mathsf D(\tau,\nu)$. If $\nu=01$, \eqref{c79:eq:i5} leads to $\cfg0\tau\nu{11}\in\mathsf C(\tau,01)$, treated below. The configurations $\cfg3\tau\nu{10}$ are the mirror images of these.
\item $\mathsf C(\tau,\nu)$. By \eqref{c79:eq:i3} its three configurations reach each other. If $y_0=0$, \eqref{c79:eq:i4} leads to $\cfg2\tau{\nu^x}{10}\in\mathsf D(\tau,\nu^x)$. If $y_{-1}=0$, the mirror image of \eqref{c79:eq:i4} leads from $\cfg2\tau\nu{11}$ to $\cfg0\tau{\nu^y}{01}\in\mathsf A(\tau,\nu^y)$. If $\nu=11$, \eqref{c79:eq:i4} and the mirror image of \eqref{c79:eq:i6} lead to $\cfg5\tau\nu{01}$ and to $\cfg4\tau\nu{01}\in\mathsf D(\tau,\nu)$.
\end{itemize}

\emph{Step 3: every irregular configuration outside $\mathcal E$ is reached from a configuration in a sector.} Call a configuration \emph{good} if it is reached from a configuration in a sector; configurations in sectors are good, and a configuration reached from a good one is good.
\begin{itemize}
\item $\cfg4\tau\nu{00}$ is good. If $y_{-1}=1$, by \eqref{c79:eq:j8} from $\cfg5\tau\nu{00}\in\mathsf A(\tau,\nu)$; if $y_0=1$, by the mirror image; and if $\nu=00$, by \eqref{c79:eq:j10} and \eqref{c79:eq:j9} from $\cfg0\tau{10}{01}\in\mathsf A(\tau,10)$, through $\cfg5\tau{10}{01}$.
\item $\cfg5\tau\nu{01}$ with $y_0=0$ is good, by \eqref{c79:eq:j10}; and so is its mirror image, $\cfg3\tau\nu{10}$ with $y_{-1}=0$.
\item If $y_0=1$, then $\cfg0\tau\nu{11}\mv2\cfg5\tau\nu{01}\mv1\cfg0\tau\nu{11}$ by \eqref{c79:eq:i4} and \eqref{c79:eq:i5}; with \eqref{c79:eq:i3}, the four configurations of $\mathsf C(\tau,\nu)\cup\{\cfg5\tau\nu{01}\}$ reach each other. By the mirror image, if $y_{-1}=1$, the four configurations of $\mathsf C(\tau,\nu)\cup\{\cfg3\tau\nu{10}\}$ reach each other. So it remains to show that $\mathsf C(\tau,\nu)$ contains a good configuration for every $(\tau,\nu)\ne(00,11)$; the remaining irregular configurations, those of $\mathsf C(00,11)$, $\cfg5{00}{11}{01}$ and $\cfg3{00}{11}{10}$, form $\mathcal E$.
\item $\nu=11$ and $w_0=1$, or $\nu=01$ and $w_0=0$. The configuration $\cfg4\tau{01}{00}$ is good, and by \eqref{c79:eq:i7} and \eqref{c79:eq:i5} it reaches $\cfg5\tau{\nu}{01}$ and $\cfg0\tau{\nu}{11}$ with $\nu=(01)^x$, which is $11$ if $w_0=1$ and $01$ if $w_0=0$.
\item $\nu=10$ and $w_{-1}=1$. By \eqref{c79:eq:j11}, from $\cfg5\tau{00}{00}\in\mathsf A(\tau,00)$.
\item The mirror images of the last two items are: $\nu=11$ and $w_{-1}=1$; $\nu=10$ and $w_{-1}=0$; $\nu=01$ and $w_0=1$. So $\mathsf C(\tau,\nu)$ contains a good configuration whenever $\nu\in\{01,10\}$, and whenever $\nu=11$ and $\tau\ne00$; and then, by the third item, $\cfg5\tau\nu{01}$ is good if $y_0=1$.
\item $\nu=00$. Let $\tau^\circ=(u,w_0)$ with $(u,w_0)\ne(0,0)$. Then $\cfg5{\tau^\circ}{\mu}{01}$ with $\mu=(\neg w_0,1)$ is good, by the last item, since $(\tau^\circ,\mu)\ne(00,11)$; and by \eqref{c79:eq:j12} it reaches $\cfg1{\tau'}{00}{11}\in\mathsf C(\tau',00)$ with $\tau'=(u\oplus[a_0w_0\ne01],w_0)$. For $w_0=1$ and $u\in\{0,1\}$ these are the cells $\tau'=01$ and $\tau'=11$; for $w_0=0$ and $u=1$, $\tau'=00$. The mirror image gives $\tau'=10$.
\end{itemize}

\emph{Step 4: from sector to sector.} For sectors write $\Sigma\to\Sigma'$ if a configuration of $\Sigma$ reaches a configuration of $\Sigma'$; then, by Step~1, every configuration of $\Sigma$ reaches every configuration of $\Sigma'$, and $\to$ is transitive. We write $\tau.\mu$ for $\Sigma(\tau,\mu)$. We claim
\begin{align}
&\cfg0\tau\nu{01}\mv{2}\cfg2\tau{\nu^x}{10}\ \text{ if }y_0=1;\label{c79:eq:k13}\\
&\cfg3\tau\mu{11}\mv{3}\cfg4{\tau'}{\mu'}{01},\label{c79:eq:k14}
\end{align}
where, for $\tau=(w_{-1},w_0)$ and $\mu=(y_{-1},y_0)$: if $w_0=0$, then $(\tau',\mu')=(\tau,\mu)$ if $w_{-1}=y_{-1}\ne y_0$, and $\tau'=(\neg w_{-1},0)$, $\mu'=(\neg y_{-1},\neg y_0)$ otherwise; if $w_0=1$, then $(\tau',\mu')=(\tau,\mu)$ if $y_{-1}=1$ and $w_{-1}=\neg(a_0\wedge y_0)$, and $\tau'=(w_{-1}\oplus a_0,1)$, $\mu'=(y_{-1},\neg y_0)$ otherwise. Indeed, \eqref{c79:eq:k13} is the row $(2,0)$ with $z_1=0\oplus(1\wedge\neg y_0)=0$. In \eqref{c79:eq:k14}, the row $(3,3)$ with $z_{-1}=1$ has $\pi=1$, so $z'_1=1\ne z_1$ and $D_r=01$, and by Lemma~\ref{c79:lem:M3}(b) and \eqref{c79:eq:br}, with $y_1=y_{-1}\oplus(\neg w_0\wedge y_0)$, $w_1=w_{-1}\oplus f(a_0,w_0,y_0)$ and $\delta=[w_1y_1\ne01]$,
\[
y'_0=y_0\oplus\delta,\qquad y'_{-1}=y_{-1}\oplus(\delta\wedge\neg w_0),\qquad w'_{-1}=w_{-1}\oplus\bigl(\delta\wedge[a_0w_0\ne01]\bigr).
\]
If $w_0=0$, then $y_1=y_{-1}\oplus y_0$ and $w_1=w_{-1}\oplus\neg y_0$, so $\delta=0$ exactly when $w_{-1}=y_{-1}=\neg y_0$, and $[a_0w_0\ne01]=1$. If $w_0=1$, then $y_1=y_{-1}$ and $w_1=w_{-1}\oplus\neg(a_0\wedge y_0)$, so $\delta=0$ exactly when $y_{-1}=1$ and $w_{-1}=\neg(a_0\wedge y_0)$, and $[a_0w_0\ne01]=a_0$.

These give four relations between sectors. In \eqref{c79:eq:k13} let $\nu=\mu^x$, where $\mu=(y_{-1},y_0)$ has $y_0=1$; then $\cfg0\tau\nu{01}\in\mathsf A(\tau,\mu^x)\subseteq\Sigma(\tau,\mu)$ and $\cfg2\tau{\mu}{10}\in\mathsf D(\tau,\mu)\subseteq\Sigma(\tau,\mu^y)$. In \eqref{c79:eq:k14}, $\cfg4{\tau'}{\mu'}{01}\in\mathsf D(\tau',\mu')\subseteq\Sigma(\tau',(\mu')^y)$, where $(\mu')^y$ is formed with $\tau'$; put $J_{a_0}(\tau,\mu)=\tau'.(\mu')^y$. So, with the mirror images,
\begin{align*}
(\mathrm S)\ \ &\tau.\mu\to\tau.\mu^y\ \text{ if }y_0=1;&(\bar{\mathrm S})\ \ &\tau.\mu\to\tau.\mu^x\ \text{ if }y_{-1}=1;\\
(\mathrm J)\ \ &\tau.\mu\to J_{a_0}(\tau,\mu);&(\bar{\mathrm J})\ \ &\tau.\mu\to\overline{J_{a_{-1}}(\bar\tau,\bar\mu)},
\end{align*}
where the bar exchanges the two bits of both cells. By the definition of $\nu^y$, $(\mu')^y=(y'_{-1},\,y'_0\oplus[w'_{-1}y'_{-1}\ne01])$ for $\tau'=(w'_{-1},w'_0)$ and $\mu'=(y'_{-1},y'_0)$: it is $\mu'$ with its second bit changed, unless $(w'_{-1},y'_{-1})=(0,1)$, and then it is $\mu'$. The following table gives, for each $\tau.\mu$, the cells $\tau'\,\mu'$ and, after the colon, the sector $J_{a_0}(\tau,\mu)=\tau'.(\mu')^y$; for $w_0=0$ they do not depend on $a_0$.
\[
\begin{array}{l|cccc}
\tau\ \backslash\ \mu&00&01&10&11\\\hline
00&10\,11\colon10.10&00\,01\colon00.00&10\,01\colon10.00&10\,00\colon10.01\\
10&00\,11\colon00.11&00\,10\colon00.10&10\,10\colon10.11&00\,00\colon00.01\\
01\ (a_0=0)&01\,01\colon01.00&01\,00\colon01.01&01\,11\colon01.11&01\,10\colon01.10\\
11\ (a_0=0)&11\,01\colon11.00&11\,00\colon11.01&11\,10\colon11.11&11\,11\colon11.10\\
01\ (a_0=1)&11\,01\colon11.00&11\,00\colon11.01&11\,11\colon11.10&01\,11\colon01.11\\
11\ (a_0=1)&01\,01\colon01.00&01\,00\colon01.01&11\,10\colon11.11&01\,10\colon01.10
\end{array}
\]
For example, for $\tau.\mu=00.00$ the exception $w_{-1}=y_{-1}\ne y_0$ does not hold, so $(\tau',\mu')=(10,11)$, and $(\mu')^y=(1,1\oplus[11\ne01])=10$; for $\tau.\mu=00.01$ it holds, so $(\tau',\mu')=(00,01)$, and $(\mu')^y=(0,1\oplus[00\ne01])=00$.

For every $c_{r-3}$, the relations $(\mathrm J)$ and $(\bar{\mathrm J})$ give the closed walk
\[
00.00\xrightarrow{\mathrm J}10.10\xrightarrow{\mathrm J}10.11\xrightarrow{\mathrm J}00.01\xrightarrow{\bar{\mathrm J}}01.00\xrightarrow{\bar{\mathrm J}}00.11\xrightarrow{\mathrm J}10.01\xrightarrow{\mathrm J}00.10\xrightarrow{\bar{\mathrm J}}00.00,
\]
which uses only the first two rows of the table: its steps $(\bar{\mathrm J})$ are the mirror images of $00.10\to10.00$, $10.00\to00.11$ and $00.01\to00.00$. Its mirror image is the closed walk
\[
00.00\xrightarrow{\bar{\mathrm J}}01.01\xrightarrow{\bar{\mathrm J}}01.11\xrightarrow{\bar{\mathrm J}}00.10\xrightarrow{\mathrm J}10.00\xrightarrow{\mathrm J}00.11\xrightarrow{\bar{\mathrm J}}01.10\xrightarrow{\bar{\mathrm J}}00.01\xrightarrow{\mathrm J}00.00 .
\]
Together they pass through all twelve sectors $\tau.\mu$ with $\tau\ne11$; so any two of these satisfy $\Sigma\to\Sigma'$. If $a_0=1$, the last two rows of the table give the closed walks
\[
01.00\xrightarrow{\mathrm J}11.00\xrightarrow{\mathrm J}01.00,\qquad01.01\xrightarrow{\mathrm J}11.01\xrightarrow{\mathrm J}01.01,\qquad01.10\xrightarrow{\mathrm J}11.10\xrightarrow{\mathrm J}11.11\xrightarrow{\mathrm J}01.10,
\]
so each of the four sectors with $\tau=11$ lies on a closed walk through a sector with $\tau=01$; if $a_{-1}=1$, the mirror images of these walks do the same with $\tau=10$. Finally let $c_{r-3}=00$ and $\tau=11$. Then
\[
\cfg3\tau{01}{00}\mv{-3}\cfg4\tau{01}{00}\mv{2}\cfg5\tau{11}{01}\mv{-3}\cfg4\tau{11}{01}
\]
by the mirror image of \eqref{c79:eq:j8}, by \eqref{c79:eq:i7} with $(01)^x=11$, and by the mirror image of \eqref{c79:eq:i6}; the first configuration lies in $\mathsf D(\tau,01)\subseteq\Sigma(\tau,00)$ and the last in $\mathsf D(\tau,11)\subseteq\Sigma(\tau,10)$, as $(01)^y=00$ and $(11)^y=10$ for $\tau=11$. With the table, $(\bar{\mathrm S})$ and $(\mathrm S)$ this gives the closed walk
\[
11.00\to11.10\xrightarrow{\mathrm J}11.11\xrightarrow{\bar{\mathrm S}}11.01\xrightarrow{\mathrm S}11.00 .
\]
To sum up: if $c_{r-3}\ne00$, any two of the sixteen sectors satisfy $\Sigma\to\Sigma'$; if $c_{r-3}=00$, this holds for any two of the twelve sectors with $\tau\ne11$, and for any two of the four sectors with $\tau=11$.

\emph{Conclusion.} Let $c'$ be as in the proposition. By Step~2, $c$ lies in a sector or reaches a configuration in one; let this sector be $\tau_1.\mu_1$. By Step~3, $c'$ lies in a sector or is reached from a configuration in one; let this sector be $\tau_2.\mu_2$. If $c_{r-3}\ne00$, then $\tau_1.\mu_1\to\tau_2.\mu_2$ by Step~4, and $c$ reaches $c'$. Let $c_{r-3}=00$. By Lemma~\ref{c79:lem:M3}(b), a move at a time $0\le m\le3$ does not change $w_0$, and it changes $w_{-1}$ only if $[a_0w_0\ne01]=1$, that is, here, only if $w_0=0$; by the mirror image, a move at a time $-3\le m\le0$ does not change $w_{-1}$, and changes $w_0$ only if $w_{-1}=0$. So no move at a time $|m|\le3$ changes whether the cell $r-2$ is $11$. Hence $\tau_1=11$ if and only if $c_{r-2}=11$, and $\tau_2=11$ if and only if $c'_{r-2}=11$; by the assumption on $c'$, $\tau_1=11$ if and only if $\tau_2=11$. So again $\tau_1.\mu_1\to\tau_2.\mu_2$ by Step~4, and $c$ reaches $c'$.
\end{proof}

The exceptions are genuine: \emph{if $c\in\Omega_r$, $|m|\le3$ and $T_mc\in\mathcal E$, then $c\in\mathcal E$}; so $\mathcal O_3(c)\cap\mathcal E=\emptyset$ for $c\notin\mathcal E$. Since $\mathcal E$ is its own mirror image, it suffices to take $0\le m\le3$ and $D=T_mc\ne c$ in $\mathcal E$, and to go through the table of Lemma~\ref{c79:lem:M3}; recall that $D_{r-2}=00$, $D_{r-1}=11$, and that $(\vartheta',D_r)$ is $(0,11)$, $(1,11)$, $(2,11)$, $(3,10)$ or $(5,01)$. The rows with $\vartheta'=4$ do not occur; nor do those with $(\vartheta',D_r)=(3,00)$, $(2,10)$, $(1,10)$, nor the row $(3,5)$ with $z_2=0$, where $D_r=(\neg y'_0,1)=01$. In the rows $(1,5)$ with $z_{-1}=1$ and $(2,4)$ with $z_0=1$, $D_r=c_r$ would be $10$ with $\vartheta=5$, respectively $11$ with $\vartheta=4$, the cell of the vacuum. In the rows with $m=0$ and $\vartheta'=1$, and in the row $(1,1)$ with $z_0=1$, no cell changes, and $c$ is $\cfg2{00}{11}{11}$, $\cfg0{00}{11}{11}$ or $\cfg1{00}{11}{11}$. In the row $(1,5)$ with $z_{-1}=0$, $D_r=11$ forces $z_0=1$, and $c=\cfg5{00}{11}{01}$. In the row $(2,0)$ with $z_1=1$, $D_r=(0,z_0)=01$ and $z'_0=z_0$, so $c_{r-1}=11$ and $z_{-1}=z_1=1$, and $c=\cfg0{00}{11}{11}$. In the row $(2,4)$ with $z_0=0$, $D_r=01$ forces $z_{-1}=0$ and $y_{-1}=0$, and then $y'_{-1}=y_{-1}\oplus[w_0y_0\ne01]=[01\ne01]=0$, which is not the case. So the five configurations of $\mathcal E$ with given cells on $(-\infty,r-3]$ leave their class but cannot be reached from it, and $(\mathrm F^{93}_3)$ fails exactly by $\mathcal E$; the condition on $(1,1)$ is invariant, as was shown at the end of the proof.

\begin{theorem}\label{c79:thm:Fk}
Let $k\ge4$ and $c\in\Omega_r$. Then $\mathcal O_k(c)$ is the set of $c'\in\Omega_r$, of any phase, with
\begin{enumerate}
\item[(i)] $c'_i=c_i$ for every $i\le r-k$; and
\item[(ii)] if the exterior cell $c_{r-k}$ is $(0,0)$, then $c'_{r-k+1}=(1,1)$ if and only if $c_{r-k+1}=(1,1)$.
\end{enumerate}
In particular $c\in\mathcal O_k(c')$ for every $c'\in\mathcal O_k(c)$.
\end{theorem}

We call the statement for a given $k$ $(\mathrm F^{93}_k)$. It has the form of Theorem~\ref{cl:thm:Fk}, with every phase of the vacuum allowed.

\begin{proof}
\emph{The base case $k=4$.} Let $\mathcal C$ be the set of $b\in\Omega_r$, of any phase, with $b=c$ on $(-\infty,r-4]$, let $\mathcal S\subseteq\mathcal C$ be the set described by (i) and (ii), and write $b_{r-4+j}=(u_j,v_j)$ for $j=0,\dots,4$, so that $(u_0,v_0)=c_{r-4}$, $(u_1,v_1)$ is the tip and $b_r=(u_4,v_4)$ the edge cell. Write $b\to b'$ if $b'\in\mathcal O_4(b)$; this relation is transitive.

\emph{Invariance.} Every element of $\mathcal O_4(c)$ lies in $\mathcal S$. Property (i) is Lemma~\ref{c79:lem:Tbasic}(b). For (ii), let $u_0=v_0=0$ and consider one move at a time $|m|\le4$ applied to some $b\in\mathcal C$. By Lemma~\ref{c79:lem:Tbasic}(b) the tip changes only if $|m|=4$. By Lemma~\ref{c79:lem:onemove}(a) (with $j=3$ and $j=2$), a move at time $4$ changes neither $v_1=Y_b(0,r-3)$ nor $Y_b(1,r-3)=u_1\oplus f(0,v_1,v_2)$, and $f(0,1,c)=1$; so if $v_1=1$, then $u_1$ does not change. By Lemma~\ref{c79:lem:Tbasic}(c), a move at time $-4$ is the $\rho$-conjugate of one at time $4$; it changes neither $u_1$ nor, if $u_1=1$, $v_1$. So $u_1v_1$ is invariant.

\emph{Transfer.} By Lemma~\ref{c79:lem:Tbasic}(c), $F^{-1}T_mF=T_{m+1}$, so $\mathcal O_3(b)\subseteq\mathcal O_4(b)$, $F^{-1}\mathcal O_3(Fb)\subseteq\mathcal O_4(b)$, and $\rho\,\mathcal O_4(b)=\mathcal O_4(\rho b)$. The cells $r-4$, $r-3$, $r-2$ of $Fb\in\Omega_r$ are $(v_0,\cdot)$, $(v_1,\alpha)$ and $(v_2,\beta)$, where $\alpha=u_1\oplus f(v_0,v_1,v_2)$ and $\beta=u_2\oplus f(v_1,v_2,v_3)$. Let $\tilde b\in\Omega_r$, of any phase, agree with $Fb$ on $(-\infty,r-3]$, with $\tilde b_{r-2}=(x,x')$ and $\tilde b_{r-1}=(x'',\cdot)$. Then $b'=F^{-1}\tilde b\in\mathcal C$ by Lemma~\ref{c79:lem:glue}(b), the rows $0$ and $1$ of $Y_{b'}$ are the $x$- and $y$-rows of $\tilde b$, and \eqref{eq:rel} at $(0,r-3)$ and $(0,r-2)$ gives
\[
b'_{r-3}=\bigl(\alpha\oplus f(v_0,v_1,x),\,v_1\bigr),\qquad b'_{r-2}=\bigl(x'\oplus f(v_1,x,x''),\,x\bigr).
\]
By Proposition~\ref{c79:prop:F3} for $Fb$, whose exterior cell and tip at level $3$ are $(v_1,\alpha)$ and $(v_2,\beta)$, $b\to b'$ whenever $\tilde b\notin\mathcal E$ and, if $(v_1,\alpha)=(0,0)$, $[\tilde b_{r-2}=(1,1)]=[(v_2,\beta)=(1,1)]$.

\emph{Step 1 (same tip).} Let $b,b'\in\mathcal C$ have the same tip $t$, and if $t=(0,0)$ let $[b_{r-2}=(1,1)]=[b'_{r-2}=(1,1)]$. Then $b\to b'$. If $b'\notin\mathcal E$, this is Proposition~\ref{c79:prop:F3} for $b$, whose exterior cell at level $3$ is $t$. Let $b'\in\mathcal E$, so $(u'_2,v'_2)=(0,0)$ and $(u'_3,v'_3)=(1,1)$, and let $b''$ be $b$ with its cell $r-2$ replaced by $(1,0)$; then $b\to b''$, as $b''\notin\mathcal E$. Since $v'_2=v''_2=0$ and the tips agree, $Fb'$ and $Fb''$ agree on $(-\infty,r-3]$; their cells $r-2$ are $(0,\cdot)\ne(1,1)$, and $(Fb')_{r-1}=(1,\,1\oplus f(0,1,v'_4))=(1,0)$, so $Fb'\notin\mathcal E$. By the transfer with $\tilde b=Fb'$, $b''\to b'$.

\emph{Step 2 (tip $(u_1,0)\to(\neg u_1,0)$).} Every $b\in\mathcal S$ with $v_1=0$ has $b\to b'$ for some $b'$ with tip $(\neg u_1,0)$; if $u_1=1$, $b'$ can be chosen with $b'_{r-2}=(1,1)$, and also with $b'_{r-2}\ne(1,1)$. By Step~1 we may first change $b_{r-2}$: if $u_1=0$ and $(u_2,v_2)=(0,1)$ we replace it by $(1,0)$, and if $u_1=1$ we choose $v_2$ freely. Now $\alpha=u_1\oplus f(v_0,0,v_2)=u_1\oplus\neg v_2$. If $(v_1,\alpha)=(0,0)$, the tip $(v_2,\beta)$ of $Fb$ is not $(1,1)$: $v_2=1$ and $\beta=u_2\oplus f(0,1,v_3)=\neg u_2=1$ would mean $(u_2,v_2)=(0,1)$ and $\alpha=u_1=0$, which we excluded. Take $\tilde b$ with $\tilde b_{r-2}=(\neg v_2,v_2)$, which is neither $(0,0)$ nor $(1,1)$, and the other cells to the right of $r-3$ and the phase arbitrary. The transfer gives $b\to b'$ with $b'_{r-3}=(\alpha\oplus f(v_0,0,\neg v_2),0)=(\neg u_1,0)$ and $b'_{r-2}=(v_2\oplus f(0,\neg v_2,x''),\neg v_2)$, which is $(1,1)$ if $v_2=0$ and $(x'',0)$ if $v_2=1$.

\emph{Step 3 (tip $(0,v_1)\to(0,\neg v_1)$).} The $\rho$-image of Step~2: every $b\in\mathcal S$ with $u_1=0$ has $b\to b'$ with tip $(0,\neg v_1)$, and if $v_1=1$ with either value of $[b'_{r-2}=(1,1)]$.

\emph{Step 4.} If $v_0=1$, every $b\in\mathcal S$ with $v_1=1$ has $b\to b'$ with tip $(\neg u_1,1)$: here $\alpha=u_1\oplus f(1,1,v_2)=u_1\oplus\neg v_2$, the exterior cell $(1,\alpha)$ of $Fb$ is not $(0,0)$, and $\tilde b_{r-2}=(\neg v_2,1)$ gives $b'_{r-3}=(\alpha\oplus f(1,1,\neg v_2),1)=(\neg u_1,1)$. If $u_0=1$, the $\rho$-image gives $b\to b'$ with tip $(1,\neg v_1)$ for every $b\in\mathcal S$ with $u_1=1$.

\emph{Conclusion.} For a tip $t$ let $A_t=\{b\in\mathcal S:\ b_{r-3}=t\}$, and split $A_{00}=A^1_{00}\cup A^0_{00}$ according to whether $b_{r-2}=(1,1)$. By Step~1, $b\to b'$ for any two elements of one of $A_{01}$, $A_{10}$, $A_{11}$, $A^0_{00}$, $A^1_{00}$. By Step~2 every element of $A_{00}$ reaches $A_{10}$ and every element of $A_{10}$ reaches both $A^0_{00}$ and $A^1_{00}$; by Step~3 every element of $A_{00}$ reaches $A_{01}$ and every element of $A_{01}$ reaches both. So $b\to b'$ for all $b,b'\in A_{00}\cup A_{01}\cup A_{10}$. If $c_{r-4}=(0,0)$, $\mathcal S$ is this union or $A_{11}$. Otherwise $v_0=1$ or $u_0=1$, $\mathcal S$ also contains $A_{11}$, and Step~4 gives $b\to b'$ in both directions between $A_{11}$ and $A_{01}$ (if $v_0=1$) or $A_{10}$ (if $u_0=1$). So $b\to b'$ for all $b,b'\in\mathcal S$; with the invariance, $\mathcal O_4(c)=\mathcal S$, which is $(\mathrm F^{93}_4)$.

\emph{The step.} Let $k\ge5$ and assume $(\mathrm F^{93}_{k-1})$. Let $\mathcal C$ be the set of $b\in\Omega_r$, of any phase, with $b=c$ on $(-\infty,r-k]$. For $b\in\mathcal C$ write $b_{r-k+j}=(u_j,v_j)$ for $j=0,\dots,3$, where $(u_0,v_0)=c_{r-k}$, and $\pi(b)=(u_1,v_1,u_2,v_2,u_3,v_3)$; as $k\ge5$, the cells $r-k+1,\dots,r-k+4$ are not the edge cell. For $\ell\in\{-1,0,1\}$ let $\sim_\ell$ be the relation on $\mathcal C$ of reachability by moves at the times $|m-\ell|\le k-1$. By Lemma~\ref{c79:lem:Tbasic}(c), $b\sim_\ell b'$ if and only if $F^\ell b'\in\mathcal O_{k-1}(F^\ell b)$, so by $(\mathrm F^{93}_{k-1})$ each $\sim_\ell$ is an equivalence relation. The three windows together are $|m|\le k$, so $\mathcal O_k(c)$ is the class of $c$ under the equivalence relation generated by $\sim_{-1}$, $\sim_0$, $\sim_1$.

We claim that $b\sim_\ell b'$ if and only if $(\pi(b),\pi(b'))\in\mathcal R_\ell$, where, writing $ab$ for $a\wedge b$ and $f=f_{93}$, the relations $\mathcal R_\ell\subseteq\{0,1\}^6\times\{0,1\}^6$, which depend on $(u_0,v_0)$, are defined as follows, with $\pi(b)=(u_1,\dots,v_3)$ and $\pi(b')=(u'_1,\dots,v'_3)$:
\begin{itemize}
\item[$\mathcal R_0$:] $(u'_1,v'_1)=(u_1,v_1)$, and $(u'_2,v'_2,u'_3,v'_3)$ is arbitrary, except that $u'_2v'_2=u_2v_2$ if $(u_1,v_1)=(0,0)$.
\item[$\mathcal R_1$:] with $\alpha=u_1\oplus f(v_0,v_1,v_2)$ and $\beta=u_2\oplus f(v_1,v_2,v_3)$: $v'_1=v_1$, $u'_1=\alpha\oplus f(v_0,v_1,v'_2)$ and $u'_2=\beta'\oplus f(v_1,v'_2,v'_3)$ for some $\beta'\in\{0,1\}$, with $v'_2,v'_3,\beta',u'_3$ arbitrary, except that $v'_2\beta'=v_2\beta$ if $(v_1,\alpha)=(0,0)$.
\item[$\mathcal R_{-1}$:] with $\gamma=v_1\oplus f(u_0,u_1,u_2)$ and $\delta=v_2\oplus f(u_1,u_2,u_3)$: $u'_1=u_1$, $v'_1=\gamma\oplus f(u_0,u_1,u'_2)$ and $v'_2=\delta'\oplus f(u_1,u'_2,u'_3)$ for some $\delta'\in\{0,1\}$, with $u'_2,u'_3,\delta',v'_3$ arbitrary, except that $\delta'u'_2=\delta u_2$ if $(\gamma,u_1)=(0,0)$.
\end{itemize}
For $\ell=0$ this is $(\mathrm F^{93}_{k-1})$ applied to $b$, whose exterior cell and tip at level $k-1$ are $(u_1,v_1)$ and $(u_2,v_2)$; for $\ell=1$ it is $(\mathrm F^{93}_{k-1})$ applied to $Fb$, whose exterior cell and tip at level $k-1$ are $(v_1,\alpha)$ and $(v_2,\beta)$, translated back by $F^{-1}$ as follows; $\ell=-1$ is the same with the rows exchanged.

 The configuration $F^\ell b$ lies in $\Omega_r$, its cells are $(Y_b(\ell-1,i),Y_b(\ell,i))$, and its exterior cell and tip at level $k-1$ are those used in the definition of $\mathcal R_\ell$. For $\ell=1$, with $\alpha=u_1\oplus f(v_0,v_1,v_2)$ as in $\mathcal R_1$, $F^{-1}$ maps the set of configurations $\tilde b$, of any phase, in $\Omega_r$, that agree with $Fb$ on $(-\infty,r-k+1]$ bijectively onto the set of $b'\in\mathcal C$ with $v'_1=v_1$ and $u'_1=\alpha\oplus f(v_0,v_1,v'_2)$. Indeed, if $b'\in\mathcal C$, then $Fb'\in\Omega_r$ by Lemma~\ref{c79:lem:glue}(b), $Fb'$ agrees with $Fb$ on $(-\infty,r-k-1]$, and its cells $r-k$ and $r-k+1$ are $\bigl(v_0,\,u_0\oplus f(Y_b(0,r-k-1),v_0,v'_1)\bigr)$ and $\bigl(v'_1,\,u'_1\oplus f(v_0,v'_1,v'_2)\bigr)$, while those of $Fb$ are the same with $v_1,u_1,v_2$ in place of $v'_1,u'_1,v'_2$; so $Fb'$ agrees with $Fb$ on $(-\infty,r-k+1]$ if and only if $v'_1=v_1$ and $u'_1=\alpha\oplus f(v_0,v_1,v'_2)$. Conversely, if $\tilde b$ is such, then $b'=F^{-1}\tilde b\in\Omega_r$ by Lemma~\ref{c79:lem:glue}(b), and $b'$ agrees with $b$ on $(-\infty,r-k]$ because $F^{-1}$ has radius $1$, so $b'\in\mathcal C$. Row $0$ of $b'$ is the row of $x$-bits of $\tilde b$, and $Y_{b'}(-1,i)=Y_{b'}(1,i)\oplus f(Y_{b'}(0,i-1),Y_{b'}(0,i),Y_{b'}(0,i+1))$; at $i=r-k+2$ this gives the formula of $\mathcal R_1$ for $u'_2$, and for $i\ge r-k+3$, for given row $0$, it is a bijection between the $y$-bits of $\tilde b_i$ and the $x$-bits of $b'_i$. The phase of $\tilde b$ is free, and so is that of $b'$, and $\tilde b_r$ differs from its vacuum cell if and only if $b'_r$ does, by Lemma~\ref{c79:lem:glue}(b). So $(\mathrm F^{93}_{k-1})$ for $Fb$ says that $b\sim_1b'$ exactly when $(\pi(b),\pi(b'))\in\mathcal R_1$, the cells $r-k+4,\dots,r$ of $b'$ and its phase being free subject to $b'\in\Omega_r$. The cases $\ell=0$ and $\ell=-1$ are the same. Hence abstract paths lift to chains of moves and chains of moves project to abstract paths, and $\mathcal O_k(c)=\{c'\in\mathcal C:\ \pi(c')\approx\pi(c)\}$, where $\approx$ is the equivalence relation generated by the $\mathcal R_\ell$.

We determine the classes of $\approx$; they depend on $(u_0,v_0)$. We use $f(a,0,c)=\neg c$, $f(a,1,1)=\neg a$, $f(0,1,c)=1$ and $f(1,1,c)=\neg c$, all from \eqref{c79:eq:f}. Write states as $(u_1,v_1,u_2,v_2,u_3,v_3)$, and divide $\{0,1\}^6$ into the five blocks
\[
C_{ab}=\{(u_1,v_1)=(a,b)\}\ \ \bigl((a,b)\ne(0,0)\bigr),\quad C_{00}^{e}=\{(u_1,v_1)=(0,0),\ u_2v_2=e\}\quad(e\in\{0,1\}).
\]
By $\mathcal R_0$, any two states of the same block are related, so each block lies in one class. The following pairs are related, for every $(u_0,v_0)$ unless stated otherwise; the free bits $v'_2,\beta',\dots$ are those of the definitions of $\mathcal R_{\pm1}$.
\begin{itemize}
\item[(E1)] $C_{10}$ and both blocks $C_{00}^0$, $C_{00}^1$, by $\mathcal R_1$. Let $(u_1,v_1)=(1,0)$; then $\alpha=1\oplus f(v_0,0,v_2)=v_2$ and $u'_1=\alpha\oplus\neg v'_2$. If $v_2=1$, then $(v_1,\alpha)\ne(0,0)$, there is no exception, and $v'_2=0$ gives $(u'_1,v'_1)=(0,0)$ with $u'_2v'_2=0$. If $v_2=0$, then $(v_1,\alpha)=(0,0)$ and the exception requires $v'_2\beta'=v_2\beta=0$; take $v'_2=1$ and $\beta'=0$, so $u'_1=0$ and $u'_2=\beta'\oplus f(0,1,v'_3)=1$, a state of $C_{00}^1$.
\item[(E2)] $C_{01}$ and $C_{00}^0$, by $\mathcal R_{-1}$. Let $(u_1,v_1)=(0,1)$ and $u_2=1$; then $\gamma=1\oplus f(u_0,0,1)=1$, there is no exception, and $v'_1=1\oplus f(u_0,0,u'_2)=u'_2$; take $u'_2=0$.
\item[(E3)] If $u_0=1$: $C_{10}$ and $C_{11}$, by $\mathcal R_{-1}$. Let $(u_1,v_1)=(1,0)$; as $u_1=1$ there is no exception, and $v'_1=\gamma\oplus f(1,1,u'_2)=\neg u_2\oplus\neg u'_2$; take $u'_2=\neg u_2$.
\item[(E4)] If $v_0=1$: $C_{01}$ and $C_{11}$, by $\mathcal R_1$. Let $(u_1,v_1)=(0,1)$; as $v_1=1$ there is no exception, and $u'_1=\alpha\oplus f(1,1,v'_2)=\neg v_2\oplus\neg v'_2$; take $v'_2=\neg v_2$.
\end{itemize}
By (E1) and (E2), the four blocks other than $C_{11}$ lie in one class. If $(u_0,v_0)\ne(0,0)$, then (E3) or (E4) joins $C_{11}$ to them, and there is one class. If $(u_0,v_0)=(0,0)$, then $u_1v_1$ is invariant: with $v_0=0$ and $v_1=1$, $\mathcal R_1$ gives $u'_1=u_1\oplus f(0,1,v_2)\oplus f(0,1,v'_2)=u_1$, and with $u_0=0$ and $u_1=1$, $\mathcal R_{-1}$ gives $v'_1=v_1$; and $\mathcal R_0$ keeps $(u_1,v_1)$. So the classes are $C_{11}=\{u_1v_1=1\}$ and the union $\{u_1v_1=0\}$ of the other four blocks. With the lifting above, this gives $(\mathrm F^{93}_k)$.
\end{proof}

The base case cannot be lowered to $k=3$, but it fails there only by the set $\mathcal E$ of Proposition~\ref{c79:prop:F3}.

\begin{corollary}\label{c79:cor:ctrl}
$(\mathrm H_k)$ holds for every $k\ge1$.
\end{corollary}

\begin{proof}
By Propositions~\ref{c79:prop:H1}, \ref{c79:prop:H2} and~\ref{c79:prop:H3} it suffices to show that $(\mathrm H_k)$ implies $(\mathrm H_{k+1})$ for $k\ge3$. Let $h\in\mathcal H$, $c\in\Omega_r$ and $t\in\Z$; replacing $c$ by $F^tc$ we may take $t=0$. There is a composite $O$ of moves at the times $|m|\le k$ with $Y_{Oc}(0,r-k+1)=Y_{Oc}(0,r-k+2)=1$: take $c'\in\mathcal C$ with the phase of $c$, with $c'_{r-k+1}$ equal to $(1,1)$ if $c_{r-k}=(0,0)$ and $c_{r-k+1}=(1,1)$ and to $(0,1)$ otherwise, with $c'_{r-k+2}=(0,1)$, and with $c'_i=c_i$ for $i\ge r-k+3$; then $c'\in\mathcal O_k(c)$ by Theorem~\ref{c79:thm:Fk} if $k\ge4$, and by Proposition~\ref{c79:prop:F3} if $k=3$, since $c'_{r-1}=(0,1)$ and so $c'\notin\mathcal E$. By Lemma~\ref{c79:lem:Tbasic}(b), $Oc$ and $Ohc$ agree with $c$ and with $hc$ on $(-\infty,r-k]$, and by Lemma~\ref{c79:lem:commute}, $Ohc=hOc$. Put $Y=Y_{Oc}$ and $G'=Y_{hOc}\oplus Y$. Then $G'(0,r-k)=Y_{hc}(0,r-k)\oplus Y_c(0,r-k)$, and by $(\mathrm H_k)$ for $Oc$, $G'$ vanishes on the columns $\ge r-k+1$ at all times. The relation \eqref{eq:rel} for both diagrams at $(0,r-k+1)$ gives, as in the proof of Lemma~\ref{cl:lem:clean},
\[
0=f\bigl(Y(0,r-k)\oplus G'(0,r-k),1,1\bigr)\oplus f\bigl(Y(0,r-k),1,1\bigr),
\]
and $f(a,1,1)=\neg a$ gives $G'(0,r-k)=0$, that is $Y_{hc}(0,r-k)=Y_c(0,r-k)$.
\end{proof}

\subsection{Conclusion for the class of 79}\label{c79:sec:end}

\begin{proof}[Proof of Lemma~\ref{c79:lem:K}]
Let $h\in\mathcal H$ have local rule $\gamma$ of radius $R$, let $w=(w_{-R},\dots,w_R)\in\A^{2R+1}$, and let $c$ equal $w$ on $[-R,R]$ and $V_0$ elsewhere. If $c=V_0$, then $\gamma(w)=(hV_0)_0=(V_0)_0=w_0$. Otherwise $c\in\Omega_r$ for some $r\in[-R,R]$, and $(\mathrm H_{R+1})$ (Corollary~\ref{c79:cor:ctrl}) gives $(hc)_i=c_i$ for $i\ge r-R$, in particular for $i=0$. So $\gamma(w)=(hc)_0=w_0$ for every $w$, and $h=\mathrm{id}$.
\end{proof}

\begin{theorem}\label{thm:79rigid}
$\Aut(X_{F_f})=\langle\sh,F_f\rangle\cong\Z^2$ for $f\in\{79,93,186,242\}$, and $(a,b)\mapsto\sh^aF_f^b$ is an isomorphism.
\end{theorem}

\begin{proof}
Let $g\in\Aut(X_{F_{93}})$. By Proposition~\ref{c79:prop:offsets}, $h=F^{-3b}\sh^{-a}F^{-j}g$ fixes $K^-$ and $K^+$, so $h=\mathrm{id}$ by Lemma~\ref{c79:lem:K}, and $g=F^j\sh^aF^{3b}\in\langle\sh,F\rangle$. If $\sh^aF^b=\mathrm{id}$, then $F^b$ fixes the constants, so $b=3n$, and $\sh^aF^{3n}$ maps $K^-$ to $\sh^{a+n}K^-$ and $K^+$ to $\sh^{a-n}K^+$; as neither kink is periodic, $a+n=a-n=0$. For $79$, $186$ and $242$ use the conjugacies of Section~\ref{sec:conj}.
\end{proof}

No computer is used in the proof for the class of $79$. Proposition~\ref{c79:prop:F3} is used twice, for $(\mathrm H_4)$ in Corollary~\ref{c79:cor:ctrl} and for the base case $(\mathrm F^{93}_4)$.
